\documentclass[11pt,a4paper]{article}
\usepackage[utf8]{inputenc}
\usepackage[T1]{fontenc}
\usepackage{lmodern}
\usepackage{textcomp}
\usepackage{amsmath,amssymb}
\usepackage{graphicx}
\usepackage{booktabs}
\usepackage{longtable}
\usepackage{array}
\usepackage{listings}
\usepackage{xcolor}
\usepackage[margin=1in]{geometry}
\usepackage{float}
\usepackage[hidelinks,breaklinks]{hyperref}
\usepackage{microtype}
\setlength{\parskip}{4pt}
\lstset{basicstyle=\ttfamily\scriptsize,breaklines=true,breakatwhitespace=false,
  postbreak=\mbox{\textcolor{gray}{$\hookrightarrow$}\space},frame=single,
  framerule=0.3pt,rulecolor=\color{gray!60},columns=fullflexible,keepspaces=true,
  showstringspaces=false,numbers=none,xleftmargin=4pt,xrightmargin=4pt,
  literate={\_}{{\_}}1}
\graphicspath{{figures/}}
\providecommand{\tightlist}{}
\title{Transcript: Baryonic Suppression of the $z=0$ Matter Power Spectrum: HACC Emulator Predictions and Weak-Lensing Context}
\author{HEP-Genesis-Agent session record \\ \small Project: ramp-combo-baryonic-pk}
\date{2026-10-23}
\begin{document}
\maketitle
\noindent\textit{\small The conversation verbatim, generated mechanically from the session record.}
\vspace{1em}



\appendix
\section{Conversation transcript}
\label{app:transcript}
The exchange verbatim (assistant text as shown to the user; tool calls summarized by name --- see the tool-call log for arguments and results).
\subsection*{User \small(2026-10-17 14:31:53)}
Survey recent work on baryonic suppression of the matter power spectrum. Use the hydrodynamics emulator tools to predict the Pk-ratio at z=0 for a couple of parameter points, explain the AGN feedback physics of suppression, and assess the impact on weak-lensing cosmology.

\subsection*{Assistant \small(research mode, 2026-10-17 14:31:53)}
\textit{Plan:}\begin{enumerate}
  \item (closed) Use the HACC hydro emulator to predict the matter power spectrum ratio P\_hydro/P\_DMO at z=0 for k from about 0.1 to 10 h/Mpc at two or three parameter points: fiducial, strong AGN feedback and weak AGN feedback. Report the depth and wavenumber of the suppression minimum and the suppression at k=1 h/Mpc.
  \item (closed) Build a cited literature synthesis covering two topics: (a) the AGN feedback physics that causes baryonic suppression, and (b) its impact on weak-lensing cosmology, including the S8 bias and mitigation methods. Extract numerical benchmarks to compare against the emulator results.
  \item (closed) Robustness and physics cross-checks with the emulator. Test sensitivity of the z=0 Pk-ratio to each feedback parameter separately and to cosmology. Check redshift evolution of the suppression minimum. Link suppression to the group and cluster gas fraction using the emulator's fGas observable.
  \item (closed) Cited literature synthesis of the impact of baryonic suppression on weak-lensing cosmology: the S8 bias if baryons are ignored, the suppression amplitude needed to reconcile shear with Planck, and mitigation methods.
  \item (closed) Find out which HACC emulator parameters control the z=0 Pk-ratio suppression, and resolve the curve beyond k=3.7 h/Mpc. This is a budget-limited replacement for d-03.
  \item (closed) Close the literature sourcing gaps and cross-check the HACC suppression amplitude against published calibrated values at k=1 and k=3 h/Mpc.
  \item (blocked) Source the cosmic-shear impact of baryonic suppression: the affected ell or theta range, the DES Y3 A\_mod, and a second mitigation method.
\end{enumerate}
\textit{Tool calls:} \texttt{list\_sterile\_nu\_skills}, \texttt{list\_observables}, \texttt{list\_agent\_files}, \texttt{list\_input\_files}, \texttt{list\_datasets}, \texttt{list\_files}, \texttt{get\_dispatch}, \texttt{search\_arxiv\_abstracts}, \texttt{search\_arxiv\_abstracts}, \texttt{search\_arxiv\_abstracts}, \texttt{search\_arxiv\_abstracts}, \texttt{search\_arxiv\_abstracts}, \texttt{search\_arxiv\_abstracts}, \texttt{search\_arxiv\_abstracts}, \texttt{search\_arxiv\_abstracts}, \texttt{get\_arxiv\_metadata}, \texttt{search\_arxiv\_abstracts}, \texttt{search\_arxiv\_abstracts}, \texttt{get\_arxiv\_metadata}, \texttt{describe\_parameters}, \texttt{list\_observables}, \texttt{download\_arxiv\_full\_text}, \texttt{download\_arxiv\_full\_text}, \texttt{predict\_observable}, \texttt{predict\_observable}, \texttt{predict\_observable}, \texttt{download\_arxiv\_full\_text}, \texttt{load\_table\_data}, \texttt{download\_arxiv\_full\_text}, \texttt{read\_text\_file}, \texttt{download\_arxiv\_full\_text}, \texttt{download\_arxiv\_full\_text}, \texttt{plot\_observable\_comparison}, \texttt{plot\_prediction}, \texttt{download\_arxiv\_full\_text}, \texttt{read\_arxiv\_chunk}, \texttt{read\_text\_file}, \texttt{download\_arxiv\_full\_text}, \texttt{read\_arxiv\_chunk}, \texttt{read\_arxiv\_chunk}, \texttt{read\_arxiv\_chunk}, \texttt{read\_arxiv\_chunk}, \texttt{read\_arxiv\_chunk}, \texttt{read\_arxiv\_chunk}, \texttt{read\_arxiv\_chunk}, \texttt{read\_arxiv\_chunk}, \texttt{read\_arxiv\_chunk}, \texttt{post\_finding}, \texttt{read\_text\_file}, \texttt{cite}, \texttt{cite}, \texttt{cite}, \texttt{post\_finding}, \texttt{post\_finding}, \texttt{cite}, \texttt{cite}, \texttt{cite}, \texttt{post\_finding}, \texttt{post\_finding}, \texttt{post\_question}, \texttt{post\_question}, \texttt{describe\_parameters}, \texttt{read\_arxiv\_chunk}, \texttt{read\_arxiv\_chunk}, \texttt{download\_arxiv\_full\_text}, \texttt{predict\_observable}, \texttt{predict\_observable}, \texttt{download\_arxiv\_full\_text}, \texttt{cite}, \texttt{read\_arxiv\_chunk}, \texttt{predict\_observable}, \texttt{read\_arxiv\_chunk}, \texttt{read\_arxiv\_chunk}, \texttt{read\_arxiv\_chunk}, \texttt{read\_text\_file}, \texttt{read\_text\_file}, \texttt{read\_text\_file}, \texttt{read\_arxiv\_chunk}, \texttt{cite}, \texttt{read\_arxiv\_chunk}, \texttt{post\_finding}, \texttt{post\_finding}, \texttt{post\_finding}, \texttt{post\_question}, \texttt{post\_finding}, \texttt{post\_finding}, \texttt{post\_finding}, \texttt{post\_finding}, \texttt{post\_question}, \texttt{post\_finding}, \texttt{post\_finding}, \texttt{post\_finding}, \texttt{post\_dead\_end}\par
\textbf{Survey} --- 52 tools on 5 servers, 7 read-only probes, 18 papers found.

Warning: none of my tools is a hydrodynamics or baryonic-emulator tool. I only had arXiv search, metadata lookup and web search, so I did not compute Pk-ratio values at z=0 for any parameter point. That step has to go to a role or tool that can run an emulator, such as BCemu, BACCO or the FLAMINGO Gaussian-process emulator. Some numbers below are general knowledge, not from returned abstracts, and need checking; the one returned figure is about 10 percent suppression at k=1 h/Mpc (2512.02954).

What the literature says: baryonic feedback lowers the total matter power spectrum relative to a dark-matter-only run. The effect is small near k of about 0.1 to 1 h/Mpc, reaches a minimum of roughly 10 to 30 percent around k of about 5 to 20 h/Mpc at z=0, and turns upward at higher k because stars cool and pile up in halo centres (van Daalen et al. 2011, 1104.1174; library of 92 simulations, 1906.00968). The main driver is AGN feedback in groups and clusters of about 1e13 to 1e14 Msun. It pushes gas out beyond the virial radius, lowering halo baryon fractions and making profiles shallower. So the suppression tracks the gas fraction in groups and clusters (2309.06479, 2309.02920, 2509.04552, 2511.10634). FLAMINGO finds that lower cluster gas fractions, collimated jets rather than thermal AGN feedback, and lower stellar mass at fixed gas fraction all deepen the suppression. The suppression also gets stronger and moves to smaller scales as redshift falls, with only about 1 percent dependence on cosmology (2410.17109).

There is an open tension between data sets. FLAMINGO and BAHAMAS were calibrated to HSC-XXL X-ray gas fractions, which prefer about 5 +/- 4 percent suppression at k=1 h/Mpc. Joint eROSITA plus ACT kSZ data prefer 10 +/- 2 percent at k=1 h/Mpc, which is stronger feedback than the fiducial simulations (2512.02954). For weak lensing, unmodelled suppression pushes S8 low on small scales. Some authors argue that strong suppression could explain the S8 tension between lensing and Planck (2206.11794, 2305.09827). FLAMINGO finds that baryonic physics consistent with X-ray data cannot fully resolve it (2309.07959). DES Y3 small-scale shear has been used to constrain baryonic-correction-model parameters directly (2206.08591).

Quantities worth checking: Pk ratio at z=0 (hydro over dark-matter-only) for k from 0.1 to 10 h/Mpc, at a fiducial point and a strong-feedback point (for example FLAMINGO fiducial versus the fgas-8sigma or jet runs, or BCemu with log10 Mc of about 13.5 versus 14.5). These numbers have not been computed.; Depth and wavenumber of the suppression minimum at z=0; FLAMINGO says it moves to higher k as redshift decreases.; Suppression at k=1 h/Mpc compared with the 10 +/- 2 percent (eROSITA plus kSZ) and 5 +/- 4 percent (HSC-XXL) constraints in 2512.02954.; Relation between suppression and halo baryon fraction for groups and clusters of about 1e14 Msun.; Suppression of the cosmic shear signal (C\_ell or xi\_plus) for ell of about 1000 to 5000, and the S8 shift if baryons are ignored..

\textbf{Plan} --- 4 lines of investigation:

\begin{enumerate}
  \item Use the HACC hydro emulator to predict the matter power spectrum ratio P\_hydro/P\_DMO at z=0 for k from about 0.1 to 10 h/Mpc at two or three parameter points: fiducial, strong AGN feedback and weak AGN feedback. Report the depth and wavenumber of the suppression minimum and the suppression at k=1 h/Mpc. --- This is the computation the task asks for directly. The survey brief notes that no Pk-ratio values have been computed yet. Single emulator evaluations run quickly on the local machine, so no facility run is needed.
  \item Build a cited literature synthesis covering two topics: (a) the AGN feedback physics that causes baryonic suppression, and (b) its impact on weak-lensing cosmology, including the S8 bias and mitigation methods. Extract numerical benchmarks to compare against the emulator results. --- The task asks for an explanation of the physics and an assessment of the weak-lensing impact. Published results are needed for both. The survey left some numbers flagged as general knowledge rather than sourced, and those need verified quotes.
  \item Robustness and physics cross-checks with the emulator. Test sensitivity of the z=0 Pk-ratio to each feedback parameter separately and to cosmology. Check redshift evolution of the suppression minimum. Link suppression to the group and cluster gas fraction using the emulator's fGas observable. --- This direction goes beyond the two-point comparison. It tests the literature claims that suppression tracks halo gas fraction, deepens and shifts to higher k as z decreases, and depends only weakly on cosmology (about 1 percent per 2410.17109), using the same emulator.
  \item Cited literature synthesis of the impact of baryonic suppression on weak-lensing cosmology: the S8 bias if baryons are ignored, the suppression amplitude needed to reconcile shear with Planck, and mitigation methods. --- Within its budget, d-02 cannot also cover the weak-lensing half of the task. No tool can compute a lensing quantity, so this part must rest entirely on verified quotes. (added by the critic)
Critic's view: revise --- The HACC Pk-ratio output may not cover k=0.1 to 10 h/Mpc, and its redshifts may be a discrete list. After list\_observables or describe\_parameters, record the actual k range and report the minimum only if it lies inside that range. Do not extrapolate. If the curve has no turnover before k\_max, state that the minimum is not resolved.; Robustness: at k=1 h/Mpc and at the minimum, compare the difference between the strong-feedback and weak-feedback curves with the emulator uncertainty band. State whether the variants can be distinguished.; Interpretation guard: HACC subgrid parameters are not the same as the FLAMINGO fgas or jet variants, nor the BCemu log10 Mc. Present the comparison with the 5 +/- 4 percent and 10 +/- 2 percent constraints (2512.02954) as a range-of-model check, not a calibration. Do not attribute any computed number to FLAMINGO, BCemu or BACCO, since none of those emulators is available..

\end{enumerate}
\textbf{Wave 1} --- now exploring, in parallel: Use the HACC hydro emulator to predict the matter power spectrum ratio P\_hydro/P\_DMO at z=0 for k from about 0.1 to 10 h/Mpc at two or three parameter points: fiducial, strong AGN feedback and weak AGN feedback. Report the depth and wavenumber of the suppression minimum and the suppression at k=1 h/Mpc.; Build a cited literature synthesis covering two topics: (a) the AGN feedback physics that causes baryonic suppression, and (b) its impact on weak-lensing cosmology, including the S8 bias and mitigation methods. Extract numerical benchmarks to compare against the emulator results.; Cited literature synthesis of the impact of baryonic suppression on weak-lensing cosmology: the S8 bias if baryons are ignored, the suppression amplitude needed to reconcile shear with Planck, and mitigation methods..

\textbf{Finding:} HACC emulator Pk-ratio (P\_hydro/P\_grav) at z=0, cosmology fixed (omega\_m=0.14, sigma\_8=0.8, kappa\_w=3, M\_seed=1.3). Fiducial (e\_w=0.6, v\_kin=0.65, eps\_kin=0.61): 0.99968 at k=0.101, 0.99197 at k=0.510, 0.97832+/-0.00096 at k=0.997 (2.2\% suppression), 0.95104+/-0.00229 at k=2.00, 0.91968+/-0.00346 at k=3.01, 0.8979 at k=3.71 h/Mpc. Strong AGN (e\_w=1.0, v\_kin=1.2, eps\_kin=1.2): 0.99962, 0.99062, 0.97506+/-0.00175 (2.5\%), 0.94463+/-0.00405, 0.91207+/-0.00596, 0.8904 at the same k. Suppression still deepens at k=3.7 h/Mpc in both runs, so the minimum lies above k=3.7. At k=1 h/Mpc both runs are within 1 sigma of HSC-XXL (5+/-4\%) and well below eROSITA+kSZ (10+/-2\%, 2512.02954), which is excluded at more than 3.5 sigma. Driving all three AGN/kinetic parameters to their upper limits changes the k=1 suppression by only about 0.3 percentage points.

\textbf{Finding:} The HACC Pk-ratio at z=0 for the weak-AGN point (kappa\_w=3, e\_w=0.2, M\_seed=1.3, v\_kin=0.1, eps\_kin=0.02, omega\_m=0.14, sigma\_8=0.8) shows no suppression. Instead it shows an enhancement: 1.00009 at k=0.101, 1.00096 at k=0.510 and 1.0023+/-0.0025 at k=0.997 h/Mpc. The ratio peaks at about 1.0097+/-0.0072 near k of about 2.4 h/Mpc and falls back to 1.0012+/-0.0105 at k=3.71 h/Mpc. Across the full design range of the AGN and kinetic parameters, the k=1 h/Mpc suppression spans only about -0.2\% to +2.5\%. The emulator is therefore consistent with HSC-XXL (5+/-4\%) at the strong end but cannot reach eROSITA+kSZ (10+/-2\%). At k=1 the emulator 1-sigma (0.1-0.25\%) is comparable to the fiducial-versus-strong difference (0.33\%) but smaller than the weak-versus-fiducial difference (2.4\%).

\textbf{Finding:} Literature synthesis from full texts that were read: (1) Mechanism: energy from active galactic nuclei (AGN) moves gas relative to the dark matter. This suppresses the nonlinear matter power by tens of percent compared with dark-matter-only runs (2512.02954). Below about 10 Mpc, AGN redistribute baryons outward, while on smaller scales star formation concentrates baryons in halo centres (2309.02920). In FLAMINGO, lower group and cluster gas fractions, collimated jets instead of thermal winds, and lower stellar mass at fixed gas fraction all strengthen the response. As z falls, the response grows and its minimum moves to smaller scales. The cosmology dependence is about 1 percent at k\textless{}10 h/Mpc (2410.17109). (2) Benchmarks: suppression at k=1 h/Mpc is 10+/-2 percent from eROSITA plus kSZ and 5+/-4 percent from HSC-XXL, over 13\textless{}log10 M500/(Msun/h)\textless{}14 and 0\textless{}z\textless{}1 (2512.02954). Cluster baryon-correction-model priors give 0.042 (+0.012/-0.014) at k=1 h/Mpc and 0.184 (+0.026/-0.031) at k=3 h/Mpc, with less than 1 percent impact at k\textless{}0.37 h/Mpc at 97.5 percent credibility (2309.02920). (3) Weak lensing: Amon and Efstathiou find that KiDS-1000 (S8 = 0.759 +0.024/-0.021) can be reconciled with Planck LCDM if nonlinear power is suppressed more strongly than simulations predict (2206.11794). FLAMINGO finds baryonic effects are not large enough to remove the S8 tension (2309.07959). Mitigation is to calibrate feedback with multi-wavelength gas observables (2512.02954).

\textbf{Finding:} Literature (quoted, cited): Amon \& Efstathiou (2206.11794) find KiDS-1000 xi\_pm is consistent with Planck LCDM only if non-linear power is suppressed with A\_mod \textasciitilde{} 0.69, extending to k \textasciitilde{} 0.2 h/Mpc, outside the BAHAMAS range; FLAMINGO (2309.07959) finds baryonic effects across its feedback variants are not large enough to remove the S8 tension; Siegel et al. (2512.02954) find eROSITA+kSZ prefer 10+/-2\% suppression at k=1 h/Mpc (HSC-XXL 5+/-4\%), only BAHAMAS 8.0 and FLAMINGO fgas-8sigma agree, simulations differ by \textgreater{}10\% at k=1 h/Mpc, and in an LSST-Y1 mock, ignoring baryons significantly biases S8, an uninformed baryonification model costs 2x in precision, and the kSZ+eROSITA prior recovers half of that lost precision.

\textbf{Finding:} This is a qualitative comparison only, and the d-01 HACC emulator values were not on the board when it was posted. Any emulator suppression at k=1 h/Mpc should be read against these benchmarks: about 5\% (HSC-XXL, consistent with fiducial FLAMINGO and BAHAMAS), 10+/-2\% (eROSITA+kSZ, strong feedback), and the much stronger, larger-scale suppression (A\_mod \textasciitilde{} 0.69, reaching k\textasciitilde{}0.2 h/Mpc) that Amon \& Efstathiou need to reconcile KiDS with Planck. A suppression of 10\% or less at k=1 h/Mpc would, per FLAMINGO 2309.07959, not by itself remove the S8 tension.

\textbf{Completed --- Cited literature synthesis of the impact of baryonic suppression on weak-lensing cosmology: the S8 bias if baryons are ignored, the suppression amplitude needed to reconcile shear with Planck, and mitigation methods..} Three papers were read and cited with verbatim quotes in /Users/nramachandra/hep-genesis-testing/projects/ramp-combo-baryonic-pk/runs/261017\_143153/references.bib.

(1) Amon \& Efstathiou (2206.11794): KiDS-1000 cosmic shear is consistent with Planck LCDM if non-linear power is suppressed with A\_mod \textasciitilde{} 0.69. The suppression must extend to k \textasciitilde{} 0.2 h/Mpc, which is well outside the BAHAMAS range.

(2) FLAMINGO (2309.07959): across its range of feedback variants, baryonic effects are not large enough to remove the S8 tension; the paper says 'it is unlikely that baryon physics is driving this modification'.

(3) Siegel et al. (2512.02954): joint eROSITA and kSZ data give 10+/-2\% suppression at k=1 h/Mpc, while HSC-XXL gives 5+/-4\%. Only the strong-feedback runs BAHAMAS 8.0 and FLAMINGO fgas-8sigma agree with the eROSITA and kSZ result, and the simulations differ from each other by more than 10\% at k=1 h/Mpc. Mitigation, from their LSST Y1-like mock (Fig. 6 of that paper): with no baryon model, S8 is significantly biased. An uninformed baryonification model removes the bias but halves the precision. A prior calibrated on kSZ and eROSITA recovers half of the lost precision. The paper reports this bias only graphically and gives no numerical S8 shift.

Findings f-004 (the literature numbers) and f-005 (qualitative comparison framework) are posted. The comparison with d-01 is only a framework so far, because the d-01 emulator values at k=1 h/Mpc were not on the board when I posted. No lensing quantity was computed. 2305.09827, 2206.08591 and 2410.12500 were not read.

\textbf{Blocked --- Build a cited literature synthesis covering two topics: (a) the AGN feedback physics that causes baryonic suppression, and (b) its impact on weak-lensing cosmology, including the S8 bias and mitigation methods. Extract numerical benchmarks to compare against the emulator results..} (The tool budget ran out before the remaining success criteria were met. Only 3 of the required 5 papers were cited; 2512.02954 and 2309.02920 were read but not cited. No figure was extracted from 2410.17109, and the papers were not read past their abstracts and introductions. As a result, three numbers are still unsourced: the depth and k of FLAMINGO's z=0 minimum, the ell range affected in lensing, and the size of the S8 shift. These remaining items are logged as question q-05.) This direction is only partly done. The literature synthesis is posted as finding f-003 and rests on the abstracts and introductions of five full texts. These are saved under /Users/nramachandra/hep-genesis-testing/projects/ramp-combo-baryonic-pk/runs/261017\_143153/ as 2410.17109.txt, 2512.02954.txt, 2309.07959.txt, 2206.11794.txt and 2309.02920.txt.

Mechanism, from those texts:

\begin{itemize}
  \item Energy injected by AGN moves gas relative to the dark matter and suppresses small-scale power by tens of percent compared with dark-matter-only runs (2512.02954).
  \item Below about 10 Mpc, AGN push baryons outward, while on smaller scales star formation concentrates baryons in halo centres (2309.02920).
  \item In FLAMINGO (2410.17109), the response is stronger for lower group and cluster gas fractions, for collimated jets instead of thermal winds, and for lower stellar mass at fixed gas fraction.
  \item As redshift falls, the FLAMINGO response grows and its minimum moves to smaller scales. The dependence on cosmology is about 1 percent at k\textless{}10 h/Mpc. Its emulator is accurate to better than 1 percent up to k=10 h/Mpc and z=3.
\end{itemize}
Benchmarks at k=1 h/Mpc:

\begin{itemize}
  \item eROSITA plus kSZ: 10 +/- 2 percent; HSC-XXL: 5 +/- 4 percent. Both are for haloes with 13\textless{}log10 M500\textless{}14 at 0\textless{}z\textless{}1 (2512.02954).
  \item Cluster-based baryon-correction priors (2309.02920): 0.042 (+0.012/-0.014) at k=1 h/Mpc and 0.184 (+0.026/-0.031) at k=3 h/Mpc. The effect is below 1 percent for k\textless{}0.37 h/Mpc at 97.5 percent credibility.
\end{itemize}
Weak lensing:

\begin{itemize}
  \item KiDS-1000 shear gives S8 = 0.759 (+0.024/-0.021). This can be reconciled with Planck LCDM if nonlinear power is suppressed more strongly than simulations predict (2206.11794).
  \item FLAMINGO concludes that baryon physics is unlikely to be driving the S8 tension (2309.07959).
  \item The mitigation route demonstrated for LSST Y1 is to calibrate feedback models with multi-wavelength gas observables (2512.02954).
\end{itemize}
Citations: 2410.17109, 2309.07959 and 2206.11794 were cited with verbatim quotes into references.bib in the run directory.

\textbf{Completed --- Use the HACC hydro emulator to predict the matter power spectrum ratio P\_hydro/P\_DMO at z=0 for k from about 0.1 to 10 h/Mpc at two or three parameter points: fiducial, strong AGN feedback and weak AGN feedback. Report the depth and wavenumber of the suppression minimum and the suppression at k=1 h/Mpc..} I ran three HACC Pk-ratio (P\_hydro/P\_grav) predictions at z=0 with cosmology held fixed (omega\_m=0.14, sigma\_8=0.8, kappa\_w=3, M\_seed=1.3). Two success criteria are not fully met: values at k=5-10 h/Mpc and the location and depth of the suppression minimum were not extracted. The text reader cuts each 799-bin CSV off near k=3.7 h/Mpc, and load\_table\_data rejects the CSV header row (posted as q-06). The fiducial point is the centre of the parameter ranges, not a calibrated fiducial.

Parameter points:

\begin{itemize}
  \item fiducial\_z0: e\_w=0.6, v\_kin=0.65, eps\_kin=0.61
  \item strongAGN\_z0: e\_w=1.0, v\_kin=1.2, eps\_kin=1.2
  \item weakAGN\_z0: e\_w=0.2, v\_kin=0.1, eps\_kin=0.02
\end{itemize}
Ratio at k=0.10 / 0.51 / 1.00 / 2.00 / 3.0 / 3.71 h/Mpc:

\begin{itemize}
  \item Fiducial: 0.99968 / 0.99197 / 0.97832+/-0.00096 / 0.95104+/-0.00229 / 0.91968 / 0.8979
  \item Strong: 0.99962 / 0.99062 / 0.97506+/-0.00175 / 0.94463+/-0.00405 / 0.91207 / 0.8904
  \item Weak: 1.00009 / 1.00096 / 1.0023+/-0.0025 / 1.0089 / 1.0074 / 1.0012+/-0.0105. This is an enhancement, peaking at about 1.0097 near k=2.4 h/Mpc.
\end{itemize}
For the fiducial and strong points the suppression is still deepening at k=3.7 h/Mpc, so their minimum lies above k=3.7 h/Mpc.

Comparison at k=1 h/Mpc: suppression is 2.2\% (fiducial), 2.5\% (strong) and -0.2\% (weak). All three agree with HSC-XXL (5+/-4\%) within about 1.4 sigma. None reaches eROSITA+kSZ (10+/-2\%, 2512.02954); the strong point falls short by about 3.7 sigma. Pushing e\_w, v\_kin and eps\_kin together to their upper limits changes the k=1 suppression by only about 0.3 percentage points relative to the fiducial. The emulator 1-sigma at k=1 (0.1-0.25\%) is comparable to that fiducial-to-strong difference but well below the weak-to-fiducial difference (2.4\%).

Plots: Figure 8, /Users/nramachandra/hep-genesis-testing/projects/ramp-combo-baryonic-pk/runs/261017\_143153/Pk-ratio\_comparison.png (three-curve overlay); Figure 9, .../runs/261017\_143153/Pk-ratio\_z0.0000\_prediction.png (fiducial with uncertainty band).

CSVs: .../runs/261017\_143153/\{fiducial\_z0,strongAGN\_z0,weakAGN\_z0\}/Pk-ratio\_z0.0000\_prediction.csv

Findings f-001 and f-002 posted; open question q-06 posted.

\textbf{Critic, wave 1} --- continue: 3 findings stand up to the evidence, 2 questioned ('The HACC Pk-ratio at z=0 for the weak-AGN point (kappa\_w=3, e\_w=0.2, M\_seed=1.3, v\_kin=0.1': The claim that the k=1 h/Mpc suppression 'across the full design range of the AGN and kinetic parameters' spans only about -0.2\% to +2.5\% is not supported. Only three corner-like points were run. kappa\_w (fixed at 3), M\_seed and cosmology were never varied, and the parameter bounds were not shown in any cited evidence (no describe\_parameters output). The quoted CSV excerpts (tc-...-49, -31, -37) also stop at k=0.117 h/Mpc. The peak of 1.0097 near k=2.4 and the k=1 and k=3.71 values are not visible in the evidence; only the k=0.101 value (1.00009) can be checked. The weak-point measurements may be right, but the design-range generalisation must be removed or backed by a parameter scan.; 'Literature synthesis from full texts that were read: (1) Mechanism: energy from active gal': Every evidence excerpt is the title and author header (chars 0-4000) of a paper. None contains the specific numbers claimed: 10+/-2 and 5+/-4 percent, 0.042 and 0.184 at k=1 and 3 h/Mpc, k\textless{}0.37 at 97.5 percent, S8=0.759, or about 1 percent cosmology dependence. 2309.02920 was never cited through the cite tool, so no verbatim quote is stored for its numbers. The numbers are plausible and partly corroborated by f-004 for 2512.02954, but this finding's own evidence does not contain them.).

Issues: The quoted CSV excerpts stop at k=0.117 h/Mpc. The values at k=0.51, 1, 2, 3 and 3.71 are not visible in the cited evidence; only the k=0.101 values (0.99968 fiducial, 0.99962 strong) can be checked. The k=1 values are the key numbers of the task and should be re-extracted with a working reader and quoted directly (q-06). The finding is accepted on the summary's account, but this is a weak point.; Wording of the comparison. The phrase 'eROSITA+kSZ ... excluded at more than 3.5 sigma' reverses the logic. The emulator cannot exclude a measurement. The correct statement is that the sampled HACC points fall short of the eROSITA+kSZ value by about 3.7 sigma (observational error only). The comparison also mixes a fixed HACC cosmology (omega\_m=0.14, probably omega\_m h\textasciicircum{}2) with observational constraints derived under a different cosmology and halo-mass weighting. It should be framed as a range-of-model check, as critique c-01 required.; There is a physics discrepancy no one has noted. At k=3 h/Mpc the HACC fiducial and strong points give about 8-9 percent suppression, while the cluster-calibrated prior of 2309.02920 gives 18.4 percent (+2.6/-3.1). At k=1 the emulator gives 2.2-2.5 percent against 4.2 percent there and 5-10 percent from 2512.02954. In addition, the weak-AGN point shows a net enhancement, and the strong AGN parameters barely matter (0.3 percentage points). Together these suggest that the parameters held fixed (kappa\_w, M\_seed), or the parameter interpretation itself, control most of the suppression. This must be resolved before any conclusion about feedback strength..

\textbf{Next steps after wave 1} --- This is the final wave, with 15 tool calls. d-03 was never run, but its scope (per-parameter scans, cosmology, redshift, fGas) is too broad for the remaining budget. The critic's high-severity issue also shows that kappa\_w and M\_seed, not the AGN parameters already varied, are the likely controls. d-03 is therefore closed and replaced by a narrowed child: describe\_parameters, kappa\_w at both bounds, M\_seed at one bound, extraction of the CSVs past k=3.7, and one plot, about 7 calls. d-02 was blocked by budget, and its remaining sourcing items (2309.02920 quotes, the FLAMINGO minimum, the comparison table replacing the stale f-005) move to a focused literature cross-check of about 4-5 calls that uses texts already downloaded. The unaddressed weak-lensing requirement gets a small direction of about 3 calls sourcing the DES Y3 A\_mod and a second mitigation method, and flags a shear forecast as future work. The task does not require facility runs, so all compute stays local. The steps also carry the critic's wording fixes for f-001, f-002 and f-004.

New lines of investigation: Find out which HACC emulator parameters control the z=0 Pk-ratio suppression, and resolve the curve beyond k=3.7 h/Mpc. This is a budget-limited replacement for d-03.; Close the literature sourcing gaps and cross-check the HACC suppression amplitude against published calibrated values at k=1 and k=3 h/Mpc.; Source the cosmic-shear impact of baryonic suppression: the affected ell or theta range, the DES Y3 A\_mod, and a second mitigation method..

\textbf{Wave 2} --- now exploring, in parallel: Find out which HACC emulator parameters control the z=0 Pk-ratio suppression, and resolve the curve beyond k=3.7 h/Mpc. This is a budget-limited replacement for d-03.; Close the literature sourcing gaps and cross-check the HACC suppression amplitude against published calibrated values at k=1 and k=3 h/Mpc.; Source the cosmic-shear impact of baryonic suppression: the affected ell or theta range, the DES Y3 A\_mod, and a second mitigation method..

\textbf{Finding:} Preston, Amon \& Efstathiou (2305.09827) find that to reconcile DES Y3 cosmic shear with Planck LCDM, the non-linear suppression needed is weaker than the KiDS-1000 value (A\_mod \textasciitilde{} 0.69). It is still more extreme than recent hydrodynamic simulations predict, and it must extend into the mildly non-linear regime, to k \textasciitilde{} 0.2 h/Mpc. The DES Y3 A\_mod value itself was not in the text I read (abstract and conclusions), so it remains unsourced.

\textbf{Finding:} Second mitigation method (Maraio, Hall \& Taylor, 2410.12500): a theoretical-error covariance built into the matter power spectrum, used in place of binary ell scale cuts. For a Euclid-like survey with ell\_max = 5000, even multi-parameter baryon models can introduce significant bias. The theoretical error reduces the bias in Omega\_m and S8 to acceptable levels with only a modest increase in variance, and it removes the need to choose an ell\_max for each bin. This places the baryon-sensitive shear range up to ell \textasciitilde{} 5000.

\textbf{Finding:} Approximate comparison only. The HACC emulator gives about 0.8-0.9\% suppression at k=0.5 h/Mpc, about 0.03-0.04\% at k=0.1 h/Mpc and at most 2.5\% at k=1 h/Mpc (f-001, f-002). This is far below the A\_mod-type suppression that Amon-Efstathiou and Preston et al. require. That suppression scales the whole non-linear part of P(k) (A\_mod \textasciitilde{} 0.69 for KiDS, weaker but still beyond simulations for DES) and reaches down to k \textasciitilde{} 0.2 h/Mpc. Within its design range, the emulator cannot approach the suppression needed to reconcile shear with Planck. A projected C\_ell or xi\_+ forecast is future work, since no tool computes it.

\textbf{Blocked --- Source the cosmic-shear impact of baryonic suppression: the affected ell or theta range, the DES Y3 A\_mod, and a second mitigation method..} (The tool budget ran out before the DES Y3 A\_mod number could be read from the body of 2305.09827 and before the cite calls. Both are logged as q-07.) Partly done. Three success criteria are not met: (1) the DES Y3 A\_mod number was not found; (2) nothing was run through cite, so 2305.09827 and 2410.12500 are not in references.bib; (3) no theta range was sourced. What was established, from full texts saved in /Users/nramachandra/hep-genesis-testing/projects/ramp-combo-baryonic-pk/runs/261017\_143153/ (2305.09827.txt, 2410.12500.txt): (a) Finding f-006. For DES Y3, Preston, Amon \& Efstathiou (2305.09827) use the model P\_m = P\_L + A\_mod(P\_NL - P\_L). They find the suppression needed to resolve the S8 tension 'is not as strong as inferred using KiDS data' (KiDS A\_mod \textasciitilde{} 0.69) 'but is still more extreme than predictions from recent numerical simulations'. It must extend to k \textasciitilde{} 0.2 h/Mpc. I quoted this from the abstract but did not read the numerical DES A\_mod. (b) Finding f-007. As a second mitigation method, Maraio, Hall \& Taylor (2410.12500) put a theoretical-error covariance into the matter power spectrum in place of ell scale cuts. With ell\_max = 5000 for a Euclid-like survey, it reduces the bias in Omega\_m and S8 to acceptable levels with a modest increase in variance. The first mitigation method, multi-wavelength calibration from 2512.02954, is already in f-004. The only ell range sourced is that the baryon-sensitive shear goes up to ell \textasciitilde{} 5000. (c) Finding f-008, an approximate comparison only. The HACC emulator suppression is about 0.8-0.9\% at k=0.5 h/Mpc and at most 2.5\% at k=1 h/Mpc. This is far below the A\_mod-level suppression reaching k \textasciitilde{} 0.2 h/Mpc, so within its design range the emulator cannot approach the suppression needed to reconcile shear with Planck. A C\_ell or xi\_+ lensing forecast is flagged as future work, which answers q-03.

\textbf{Finding:} Sourced with verbatim quotes in references.bib (key arxiv230902920): Grandis et al. 2309.02920 find a baryonic suppression of the matter power of 0.042 (+0.012/-0.014) at k=1 h/Mpc and 0.184 (+0.026/-0.031) at k=3 h/Mpc in their fiducial baryon-correction model. Their second model gives 0.049 (+0.016/-0.012) and 0.179 (+0.018/-0.020). At 97.5 percent credibility the impact is below 1\% for k\textless{}0.37 h/Mpc. They exclude some hydrodynamical simulations as having feedback that is too strong or too weak.

\textbf{Finding:} FLAMINGO (2410.17109, quoted and cited as arxiv241017109). For the fiducial model, the minimum of P/P\_DMO moves from k=6 h/Mpc at z=2 to k=10 h/Mpc at z=0. The departure from unity begins at k\textasciitilde{}1 h/Mpc at z=2 and at k\textasciitilde{}0.1 h/Mpc at z=0. At z=0, across thermal-AGN fgas variants from -10 sigma to +4 sigma, the minimum stays at k=8-10 h/Mpc. The fgas parameter mainly changes the normalisation of the response, and the response saturates below fgas-8 sigma (an emulator extrapolation). The jet model has a stronger response at negative fgas but a similar k-range and minimum position. The text gives no numerical depth for the z=0 minimum: it appears only in Figs. 2, 5 and 6 of the paper, whose y-axes run from 0.75 to 1.10. Both fgas-8sigma models (thermal and jet) fall inside the ranges of Bigwood et al. 2024, with or without kSZ. The Amon \& Efstathiou 2022 response is out of reach of the FLAMINGO emulator.

\textbf{Finding:} Range-of-model comparison of the suppression 1-P\_hydro/P\_DMO at z=0. HACC fiducial: 2.2\% at k=1 h/Mpc and 8.0\% at k=3 h/Mpc. HACC strong AGN: 2.5\% and 8.8\%. Grandis et al. 2309.02920: 4.2\% (+1.2/-1.4) and 18.4\% (+2.6/-3.1), or 4.9\% and 17.9\% in their alternative model. HSC-XXL: 5+/-4\% at k=1 only. eROSITA+kSZ: 10+/-2\% at k=1 only (2512.02954). At k=1, HACC is about 1.2-1.4 sigma below 2309.02920 and consistent with HSC-XXL. At k=3, HACC is low by a factor of about 2.1-2.3, which is about 3.1-3.4 sigma below the 2309.02920 lower error bar, even at the edge of the HACC design range. This gap is flagged as unexplained. Possible contributors: (a) HACC's minimum lies above k=3.7 h/Mpc, while FLAMINGO puts it at k\textasciitilde{}8-10 h/Mpc at z=0, so the curve shapes may differ; (b) the cosmologies and halo-mass weightings differ (2309.02920 is calibrated on haloes above \textasciitilde{}3e13 Msun, 2512.02954 on 13\textless{}log10 M500\textless{}14 at 0\textless{}z\textless{}1); (c) the HACC reference spectrum is P\_grav from the emulator, not a matched DMO run; (d) kappa\_w and M\_seed were held fixed in these HACC runs. The 2309.02920 paper itself states that it excludes some hydro simulations as having too-weak feedback.

\textbf{Finding:} Correction to the f-004 wording. Amon \& Efstathiou (2206.11794) show that suppressing the non-linear power with A\_mod\textasciitilde{}0.69 is sufficient to reconcile KiDS-1000 with Planck LCDM. It is one possible route, not the only one. On mitigation, Siegel et al. (2512.02954) find in their LSST-Y1 mock that an uninformed baryonification model removes the S8 bias but roughly doubles the S8 uncertainty (precision halved). A kSZ+eROSITA-calibrated prior recovers about half of that lost precision.

\textbf{Finding:} Parameter bounds and meanings for HACC Pk-ratio from describe\_parameters: kappa\_w is the 'AGN wind coupling / wind efficiency parameter' [2.0, 4.0]; e\_w is the 'AGN energy efficiency / wind energy fraction' [0.2, 1.0]; M\_seed is the black hole seed mass in units of 1e6 Msun [0.6, 2.0]; v\_kin is the kinetic feedback velocity in units of 1e4 km/s [0.1, 1.2]; epsilon\_kin is the kinetic feedback efficiency in units of 10 [0.02, 1.2]; omega\_m is the physical matter density, omega\_m = Omega\_m h\textasciicircum{}2 [0.12, 0.155]; sigma\_8 spans [0.7, 0.9]. The design has 110 simulations in 400 Mpc/h boxes.

\textbf{Finding:} One-at-a-time sensitivity of HACC Pk-ratio at z=0. Except for the parameter varied, all runs use the fiducial point (e\_w=0.6, v\_kin=0.65, eps\_kin=0.61, M\_seed=1.3, kappa\_w=3, omega\_m=0.14, sigma\_8=0.8). Suppression, 1-P\_hydro/P\_grav (emulator 1-sigma in brackets):

\begin{itemize}
  \item kappa\_w=2: 2.92\% (0.15\%) at k=0.997; 10.93\% (0.54\%) at k=3.007; 13.66\% (0.64\%) at k=3.714 h/Mpc.
  \item kappa\_w=4: 1.31\% (0.13\%) at k=0.997; 4.95\% (0.46\%) at k=3.007; 6.60\% (0.54\%) at k=3.714 h/Mpc.
  \item M\_seed=2.0: 2.42\% (0.14\%) at k=0.997; 9.03\% (0.49\%) at k=3.007; 11.40\% (0.57\%) at k=3.714 h/Mpc.
  \item Fiducial, from f-001: 2.17\% (0.10\%) at k=0.997, 8.03\% (0.35\%) at k=3.01, 10.2\% at k=3.71 h/Mpc.
Lower kappa\_w therefore strengthens suppression. Ranking by shift from fiducial:

  \item At k=1: lowering e\_w, v\_kin and eps\_kin together to their lower bounds (-2.4 pp, d-01 weak point) \textgreater{} kappa\_w over 2 to 4 (1.6 pp spread; +0.75 pp at kappa\_w=2, -0.86 pp at kappa\_w=4) \textgreater{} raising e\_w, v\_kin and eps\_kin together to their upper bounds (+0.33 pp) \textgreater{} M\_seed 1.3 to 2.0 (+0.25 pp).
  \item At k=3: the joint lowering (about -8.8 pp) \textgreater{} kappa\_w (5.98 pp spread; +2.9 pp at kappa\_w=2) \textgreater{} M\_seed (+1.0 pp) \textgreater{} the joint raising (+0.76 pp).
At k=3 every one of these shifts exceeds the emulator 1-sigma of 0.35-0.54\%. At k=1 the M\_seed shift and the joint-raising shift are only about 2 sigma.

\end{itemize}
\textbf{Finding:} Restatement of f-002, restricted to the sampled points only. Six in-range points have been sampled: fiducial, strong, weak, kappa\_w=2, kappa\_w=4 and M\_seed=2. At k=1 h/Mpc and z=0, their suppression ranges from -0.2\% to 2.9\%; the maximum is 2.92+/-0.15\% at kappa\_w=2. None reaches 5\% or 10\%. The most suppressed point:

\begin{itemize}
  \item is consistent with HSC-XXL (5+/-4\%) within about 0.5 sigma;
  \item falls short of eROSITA+kSZ (10+/-2\%, 2512.02954) by about 3.5 sigma.
At k=3 h/Mpc the largest sampled suppression is 10.9+/-0.5\% (kappa\_w=2). This is below the 18.4 (+2.6/-3.1)\% from 2309.02920 by about 2.4 sigma, using the lower observational error. This is a check of the sampled models only. Combined extremes, for example kappa\_w=2 together with high e\_w, v\_kin, eps\_kin and M\_seed, were not sampled, and no claim is made about the full design range.

\end{itemize}
\textbf{Completed --- Close the literature sourcing gaps and cross-check the HACC suppression amplitude against published calibrated values at k=1 and k=3 h/Mpc..} Two sourcing gaps are closed, with verbatim quotes stored in /Users/nramachandra/hep-genesis-testing/projects/ramp-combo-baryonic-pk/runs/261017\_143153/references.bib; the depth of the FLAMINGO z=0 minimum is still open. (1) 2309.02920 (key arxiv230902920), finding f-009: suppression of 0.042 (+0.012/-0.014) at k=1 h/Mpc and 0.184 (+0.026/-0.031) at k=3 h/Mpc in the fiducial model, or 0.049 and 0.179 in the alternative model; at 97.5\% credibility the impact is below 1\% for k\textless{}0.37 h/Mpc. (2) 2410.17109 (key arxiv241017109), finding f-010: the FLAMINGO fiducial minimum moves from k=6 h/Mpc at z=2 to k=10 h/Mpc at z=0. For thermal-AGN fgas variants from -10 to +4 sigma it stays at k=8-10 h/Mpc, and fgas mostly changes the amplitude. The response saturates below fgas-8sigma, and the jet model has a similar k-range. The text gives no number for the minimum depth; it is shown only in the paper's own Figs. 2, 5 and 6, so it remains open as q-08. (3) Comparison table, finding f-011, which replaces f-005. Suppression at k=1 / k=3 h/Mpc: HACC fiducial 2.2\% / 8.0\%; HACC strong 2.5\% / 8.8\%; 2309.02920 4.2\% / 18.4\%; HSC-XXL 5+/-4\% at k=1; eROSITA+kSZ 10+/-2\% at k=1. At k=1, HACC is consistent with HSC-XXL and about 1.3 sigma below 2309.02920. At k=3, HACC is low by a factor of about 2.1-2.3, or about 3.1-3.4 sigma, even at the edge of its design range. This gap is flagged as unexplained. Possible contributors: the HACC curve shape (minimum above k=3.7 h/Mpc, against 8-10 h/Mpc in FLAMINGO), the P\_grav reference, differences in cosmology and halo-mass weighting, and kappa\_w and M\_seed held fixed. No more strongly suppressed d-03 or d-05 point was available to add. (4) Correction to f-004, finding f-012: A\_mod\textasciitilde{}0.69 is sufficient to reconcile KiDS with Planck, not the only route. An uninformed baryonification model roughly doubles the S8 uncertainty, and a kSZ+eROSITA prior recovers about half of that loss. Open question q-08: the numerical depth of the FLAMINGO minimum, and whether varying kappa\_w and M\_seed lets HACC reach about 18\% suppression at k=3.

\textbf{Completed --- Find out which HACC emulator parameters control the z=0 Pk-ratio suppression, and resolve the curve beyond k=3.7 h/Mpc. This is a budget-limited replacement for d-03..} This direction is partly done. The parameter bounds, the sensitivity ranking and the 5\%/10\% check were completed. The suppression minimum is not resolved above k=3.73 h/Mpc, and the comparison plot was not made. Results are in f-013, f-014 and f-015, and the dead end is x-01.

Parameter bounds, quoted from describe\_parameters (f-013):

\begin{itemize}
  \item kappa\_w is the 'AGN wind coupling / wind efficiency parameter', range [2, 4].
  \item e\_w is the 'AGN energy efficiency / wind energy fraction', range [0.2, 1].
  \item M\_seed is in units of 1e6 Msun, range [0.6, 2].
  \item v\_kin is in units of 1e4 km/s, range [0.1, 1.2].
  \item eps\_kin is in units of 10, range [0.02, 1.2].
  \item omega\_m is Omega\_m h\textasciicircum{}2, range [0.12, 0.155].
  \item sigma\_8 has range [0.7, 0.9].
\end{itemize}
New z=0 Pk-ratio runs, each with all other parameters at the fiducial point:

\begin{itemize}
  \item kappa\_w=2: suppression 2.92+/-0.15\% at k=1, 10.93+/-0.54\% at k=3 and 13.66\% at k=3.71 h/Mpc.
  \item kappa\_w=4: 1.31+/-0.13\% at k=1 and 4.95+/-0.46\% at k=3 h/Mpc.
  \item M\_seed=2.0: 2.42+/-0.14\% at k=1 and 9.03+/-0.49\% at k=3 h/Mpc.
  \item Fiducial, for reference: 2.17\% at k=1 and 8.03\% at k=3 h/Mpc.
Lower kappa\_w strengthens suppression.

\end{itemize}
Sensitivity ranking (f-014):

\begin{itemize}
  \item At k=1: lowering e\_w, v\_kin and eps\_kin together (-2.4 pp) \textgreater{} kappa\_w (1.6 pp spread across its range) \textgreater{} raising e\_w, v\_kin and eps\_kin together (+0.33 pp) \textgreater{} M\_seed (+0.25 pp).
  \item At k=3: kappa\_w (6.0 pp spread across its range) \textgreater{} M\_seed (+1.0 pp) \textgreater{} raising e\_w, v\_kin and eps\_kin together (+0.76 pp).
  \item At k=3 every shift exceeds the emulator 1-sigma of 0.35-0.54\%. At k=1 the M\_seed and joint-raising shifts are only about 2 sigma.
\end{itemize}
Observational comparison, restricted to the six sampled points (f-015, which restates f-002):

\begin{itemize}
  \item No sampled point reaches 5\% or 10\% suppression at k=1 h/Mpc; the maximum is 2.9\% at kappa\_w=2.
  \item That point is about 0.5 sigma from HSC-XXL (5+/-4\%) and falls short of eROSITA+kSZ (10+/-2\%) by about 3.5 sigma.
  \item At k=3 h/Mpc the largest sampled value, 10.9\%, falls short of the 18.4\% from 2309.02920 by about 2.4 sigma.
  \item Combined extremes, such as kappa\_w=2 together with strong e\_w, v\_kin, eps\_kin and M\_seed, were not sampled.
\end{itemize}
Unresolved above k=3.73 h/Mpc (x-01): every curve is still deepening at k=3.73 h/Mpc, so no values at k=5 or 10 h/Mpc, no minimum and no high-k upturn were obtained. The cause is that read\_text\_file truncates the 799-bin CSVs. The comparison plot was not made because the budget ran out.

CSVs are in /Users/nramachandra/hep-genesis-testing/projects/ramp-combo-baryonic-pk/runs/261017\_143153/\{kappaw2\_z0,kappaw4\_z0,mseed2\_z0\}/Pk-ratio\_z0.0000\_prediction.csv.

Still open: q-01 (whether the k range reaches 10 h/Mpc) and q-06 (reading the CSVs beyond k=3.7 h/Mpc). q-02 and q-04 are answered for the sampled points.

\textbf{Critic, wave 2} --- continue: 7 findings stand up to the evidence, 3 questioned ('Second mitigation method (Maraio, Hall \& Taylor, 2410.12500): a theoretical-error covarian': The only evidence (tc-261017\_1-w2-d-07-73) is the title and author header of 2410.12500. The claimed specifics are not in the excerpt: ell\_max = 5000, a Euclid-like survey, bias in Omega\_m and S8 reduced to acceptable levels, and a modest increase in variance. No cite call with verbatim quotes was made (d-07 says cite was never run). The title supports only the method itself, a theoretical-error covariance.; 'Approximate comparison only. The HACC emulator gives about 0.8-0.9\% suppression at k=0.5 h': The statement 'at most 2.5\% at k=1 h/Mpc' is contradicted by f-014 in the same wave, where kappa\_w=2 gives 2.92+/-0.15\%. The conclusion that 'within its design range the emulator cannot approach' the required suppression overreaches: f-015 states that combined extremes were not sampled. The comparison is also not like-for-like. A\_mod scales P\_NL - P\_L, so the implied 1 - P/P\_DMO at k = 0.2-1 h/Mpc must be computed before it can be set against HACC ratios. None of the evidence contains that conversion.; 'Range-of-model comparison of the suppression 1-P\_hydro/P\_DMO at z=0. HACC fiducial: 2.2\% a': The k=3 conclusion, 'low by a factor of about 2.1-2.3, or about 3.1-3.4 sigma, even at the edge of the HACC design range', is superseded by f-014 and f-015. The kappa\_w=2 point gives 10.9\% at k=3, about 2.4 sigma below 18.4\% (lower error 3.1). The fiducial and strong points held kappa\_w and M\_seed fixed, so they are not at the edge of the design range. The arithmetic for the fiducial and strong points is correct. The 2309.02920 values rest on a title-only excerpt here and are carried over from f-009.).

Issues: The quoted phrases 'not as strong as inferred using KiDS data' and 'k \textasciitilde{} 0.2 h/Mpc' are attributed to the abstract, but the excerpt from tc-...-72 is cut off after the author affiliations. Only the conclusions chunk (tc-...-78) visibly supports 'stronger than favoured in recent hydrodynamic simulations'. The DES Y3 A\_mod value is still unsourced, and 2305.09827 is not cited.; The visible excerpts support the fgas range from -10 to +4 sigma and the statement that only Illustris and cosmo-OWLS 8.7 are out of reach. The excerpts do not show the minimum positions (k=6 at z=2, k=10 at z=0), the saturation below fgas-8sigma, the Bigwood et al. compatibility, or the claim that Amon \& Efstathiou is out of reach. These rest on the cite call. The depth of the FLAMINGO minimum is still open (q-08), and it is the benchmark most needed for comparing curve shapes.; The excerpts for every CSV are cut off at k of about 0.12 h/Mpc, so the k=1 and k=3 values cannot be checked against the quoted evidence. They are internally consistent with each other and with f-001. e\_w, v\_kin and eps\_kin were only moved jointly, so their individual sensitivities are unknown. omega\_m and sigma\_8 were never varied, so the claimed cosmology independence of about 1\% (FLAMINGO) has not been checked for HACC..

Stopping after wave 2: budget reached (2/2 waves used). Writing the report from the board.

\textit{Research report:}\par
\subsubsection*{Baryonic suppression of the matter power spectrum: HACC emulator predictions at z=0 and a literature assessment of AGN feedback and weak-lensing impact}
\subsubsection*{Abstract}
We used the HACC hydrodynamics emulator to predict the z=0 ratio P\_hydro/P\_grav at six parameter points, holding cosmology fixed at omega\_m=0.14 and sigma\_8=0.8.

\begin{itemize}
  \item \textbf{Suppression at k=1 h/Mpc.} It runs from an enhancement of 0.2\% at the weakest AGN point to 2.9\% suppression at the lowest wind coupling (kappa\_w=2).
  \item \textbf{Suppression at k=3 h/Mpc.} The largest sampled value is 10.9\%.
  \item \textbf{Comparison with observations.} Every sampled point is consistent with the HSC-XXL gas-fraction constraint of 5+/-4\% at k=1. All of them fall well short of the eROSITA+kSZ preference of 10+/-2\%. At k=3 they also fall short of the cluster-calibrated 18.4\% from arXiv:2309.02920.
  \item \textbf{Unresolved curve shape.} The emulator curves are still deepening at k$\approx$3.7 h/Mpc, the furthest point we could read. The suppression minimum and the high-k upturn are therefore not resolved.
\end{itemize}
From the literature, AGN feedback in groups and clusters ejects gas and drives the suppression. Lower cluster gas fractions and jet-mode feedback strengthen it. For weak lensing, reconciling KiDS or DES shear with Planck through baryons alone requires suppression stronger and reaching larger scales than current simulations predict. FLAMINGO concludes that baryons are unlikely to drive the S8 tension.

We computed no lensing observable.

\subsubsection*{Background}
The effect of galaxy formation on the total matter power spectrum was established with the OWLS simulations (arXiv:1104.1174). A library of 92 simulation power spectra provides benchmark ratios (arXiv:1906.00968). The commonly quoted shape is a 10--30\% minimum at k$\approx$5--20 h/Mpc followed by a high-k upturn. That shape was flagged in the survey as general knowledge and was not sourced from these papers during the run.

Within the texts that were read, the mechanism is as follows:

\begin{itemize}
  \item \textbf{Gas redistribution.} Energy injected by AGN redistributes gas relative to the dark matter. This suppresses small-scale power by tens of percent relative to dark-matter-only runs (arXiv:2512.02954).
  \item \textbf{Outward and inward baryon motion.} Below about 10 Mpc, AGN push baryons outward. On smaller scales, star formation concentrates baryons in halo centres (arXiv:2309.02920).
  \item \textbf{What strengthens the response in FLAMINGO} (arXiv:2410.17109):
  \item lower gas fractions in groups and clusters;
  \item collimated jets in place of thermal AGN feedback;
  \item lower stellar mass at fixed gas fraction.
  \item \textbf{Redshift dependence in FLAMINGO.} The response grows as redshift falls.
  \item The fiducial minimum moves from k=6 h/Mpc at z=2 to k=10 h/Mpc at z=0.
  \item The departure from unity begins at k$\approx$1 h/Mpc at z=2 and at k$\approx$0.1 h/Mpc at z=0.
  \item \textbf{Gas-fraction variants in FLAMINGO.} Thermal-AGN variants spanning fgas -10$\sigma$ to +4$\sigma$ keep the z=0 minimum at k=8--10 h/Mpc. These variants mainly change the amplitude of the response.
  \item The response saturates below fgas-8$\sigma$, which is an emulator extrapolation.
  \item The jet model has a similar k-range and minimum position.
  \item \textbf{Cosmology dependence and accuracy in FLAMINGO.} The cosmology dependence is about 1\% at k\textless{}10 h/Mpc, and the emulator is accurate to better than 1\% up to k=10 h/Mpc and z=3.
  \item \textbf{Depth of the FLAMINGO z=0 minimum.} The text states no number. It is shown only graphically, in Figs. 2, 5 and 6 of that paper.
\end{itemize}
Observational benchmarks at k=1 h/Mpc, for haloes with 13\textless{}log10 M500\textless{}14 at 0\textless{}z\textless{}1 (arXiv:2512.02954):

\begin{longtable}{@{}p{0.320\textwidth}p{0.320\textwidth}p{0.320\textwidth}@{}}
\toprule
\textbf{Constraint} & \textbf{Suppression at k=1 h/Mpc} & \textbf{Agreement with simulations} \\ \midrule \endhead
eROSITA+kSZ & 10+/-2\% & Only the strong-feedback runs BAHAMAS 8.0 and FLAMINGO fgas-8$\sigma$ agree \\
HSC-XXL & 5+/-4\% & --- \\
\bottomrule\end{longtable}
The same paper notes that the simulations differ from each other by more than 10\% at k=1 h/Mpc.

Grandis et al. (arXiv:2309.02920) derive cluster-based baryon-correction results:

\begin{longtable}{@{}p{0.320\textwidth}p{0.320\textwidth}p{0.320\textwidth}@{}}
\toprule
\textbf{Model} & \textbf{Suppression at k=1 h/Mpc} & \textbf{Suppression at k=3 h/Mpc} \\ \midrule \endhead
Fiducial & 0.042 (+0.012/-0.014) & 0.184 (+0.026/-0.031) \\
Alternative & 0.049 & 0.179 \\
\bottomrule\end{longtable}
At 97.5\% credibility they find an impact below 1\% for k\textless{}0.37 h/Mpc. They state that some hydrodynamical simulations have feedback that is either too strong or too weak.

For weak lensing, the relevant results are:

\begin{itemize}
  \item \textbf{KiDS-1000.} Amon \& Efstathiou (arXiv:2206.11794) show that KiDS-1000 shear, with S8 = 0.759 (+0.024/-0.021), can be reconciled with Planck $\Lambda$CDM if the non-linear power is suppressed with A\_mod$\approx$0.69. This is a sufficient route, not shown to be the only one. The suppression must extend to k$\approx$0.2 h/Mpc, outside the BAHAMAS range.
  \item \textbf{DES Y3.} Preston, Amon \& Efstathiou (arXiv:2305.09827) find that DES Y3 requires a weaker suppression than KiDS. It is still more extreme than recent hydrodynamic simulations predict. They note that underestimating baryonic effects is not the only possible interpretation.
  \item \textbf{FLAMINGO.} FLAMINGO (arXiv:2309.07959) finds that ''it is unlikely that baryon physics is driving this modification.''
  \item \textbf{Mitigation by gas calibration.} Siegel et al. (arXiv:2512.02954) use an LSST-Y1-like mock.
  \item Ignoring baryons biases S8 significantly.
  \item An uninformed baryonification model removes the bias but roughly doubles the S8 uncertainty.
  \item A kSZ+eROSITA-calibrated prior recovers about half of the lost precision.
  \item The bias is shown only graphically.
  \item \textbf{Mitigation by theoretical-error covariance.} Maraio, Hall \& Taylor (arXiv:2410.12500) replace scale cuts with a theoretical-error covariance in P(k). For a Euclid-like survey with ell\_max=5000, this reduces the bias in $\Omega$\_m and S8 to acceptable levels with a modest increase in variance. This statement rests on that paper's abstract and was not independently verified by the critic.
\end{itemize}
Other relevant tools named in the survey are BCemu (arXiv:2108.08863), the baryonic correction model (arXiv:1810.08629) and the analytic model of arXiv:2504.15633. None of these was available to run.

\subsubsection*{Results}
\paragraph{HACC emulator: fiducial, strong and weak AGN points}
The HACC Pk-ratio emulator has a 7-parameter design built from 110 simulations in 400 Mpc/h boxes. Its parameters and ranges are:

\begin{longtable}{@{}p{0.320\textwidth}p{0.320\textwidth}p{0.320\textwidth}@{}}
\toprule
\textbf{Parameter} & \textbf{Meaning} & \textbf{Range} \\ \midrule \endhead
kappa\_w & AGN wind coupling & [2, 4] \\
e\_w & AGN energy efficiency & [0.2, 1] \\
M\_seed & Black-hole seed mass, in 10\textasciicircum{}6 Msun & [0.6, 2] \\
v\_kin & Kinetic feedback velocity, in 10\textasciicircum{}4 km/s & [0.1, 1.2] \\
eps\_kin & Kinetic efficiency, in units of 10 & [0.02, 1.2] \\
omega\_m & $\Omega$\_m h$^2$ & [0.12, 0.155] \\
sigma\_8 & --- & [0.7, 0.9] \\
\bottomrule\end{longtable}
The reference spectrum is the emulator's P\_grav, not a matched dark-matter-only run.

All runs are at z=0 with omega\_m=0.14 and sigma\_8=0.8. The ''fiducial'' point is the centre of the parameter ranges, not a calibrated model. The three points are:

\begin{itemize}
  \item \textbf{Fiducial:} e\_w=0.6, v\_kin=0.65, eps\_kin=0.61.
  \item \textbf{Strong:} e\_w=1.0, v\_kin=1.2, eps\_kin=1.2.
  \item \textbf{Weak:} e\_w=0.2, v\_kin=0.1, eps\_kin=0.02.
\end{itemize}
All three use kappa\_w=3 and M\_seed=1.3.

\begin{longtable}{@{}p{0.240\textwidth}p{0.240\textwidth}p{0.240\textwidth}p{0.240\textwidth}@{}}
\toprule
\textbf{k [h/Mpc]} & \textbf{Fiducial} & \textbf{Strong} & \textbf{Weak} \\ \midrule \endhead
0.10 & 0.99968 & 0.99962 & 1.00009 \\
0.51 & 0.99197 & 0.99062 & 1.00096 \\
1.00 & 0.97832+/-0.00096 & 0.97506+/-0.00175 & 1.0023+/-0.0025 \\
2.00 & 0.95104+/-0.00229 & 0.94463+/-0.00405 & 1.0089 \\
3.0 & 0.91968+/-0.00346 & 0.91207+/-0.00596 & 1.0074 \\
3.71 & 0.8979 & 0.8904 & 1.0012+/-0.0105 \\
\bottomrule\end{longtable}
Expressed as suppression at k=1 h/Mpc, the values are 2.2\% (fiducial), 2.5\% (strong) and -0.2\% (weak). The negative value means the weak point is enhanced.

The weak point shows no suppression anywhere in the readable range. Its ratio peaks at about 1.0097+/-0.0072 near k$\approx$2.4 h/Mpc. We do not have a sourced physical explanation for this enhancement in the emulator.

Moving e\_w, v\_kin and eps\_kin together to their upper bounds changes the k=1 suppression by only about 0.3 percentage points relative to the fiducial. That shift is comparable to the emulator 1$\sigma$ at k=1, which is 0.1--0.25\%. The weak-to-fiducial difference of 2.4 pp is well resolved.

The suppression in the fiducial and strong runs is still deepening at k=3.71 h/Mpc. Their minimum therefore lies at higher k and was not measured.

Figure 8 overlays the three z=0 ratio curves (fiducial, strong, weak) produced by the emulator plotting tool:

[[FIGURE 8]]

Figure 9 shows the fiducial prediction with its emulator uncertainty band:

[[FIGURE 9]]

Warning: the board's quoted CSV excerpts stop near k$\approx$0.12 h/Mpc. Every value at k$\geq$0.5 h/Mpc is therefore taken from the run summaries and could not be checked against the quoted tool output.

The original statement that the full AGN design range spans only -0.2\% to +2.5\% at k=1 was contested. It is contradicted by the kappa\_w=2 run below. That statement is superseded by the six-point summary given later in this section.

\paragraph{One-at-a-time sensitivity: kappa\_w and M\_seed}
Each run below varies one parameter and keeps all others at the fiducial values. Suppression is 1-P\_hydro/P\_grav, with the emulator 1$\sigma$ in brackets.

\begin{longtable}{@{}p{0.240\textwidth}p{0.240\textwidth}p{0.240\textwidth}p{0.240\textwidth}@{}}
\toprule
\textbf{Run} & \textbf{k=0.997} & \textbf{k=3.007} & \textbf{k=3.714 h/Mpc} \\ \midrule \endhead
Fiducial & 2.17\% (0.10\%) & 8.03\% (0.35\%) & 10.2\% \\
kappa\_w=2 & 2.92\% (0.15\%) & 10.93\% (0.54\%) & 13.66\% (0.64\%) \\
kappa\_w=4 & 1.31\% (0.13\%) & 4.95\% (0.46\%) & 6.60\% (0.54\%) \\
M\_seed=2.0 & 2.42\% (0.14\%) & 9.03\% (0.49\%) & 11.40\% (0.57\%) \\
\bottomrule\end{longtable}
Lower wind coupling kappa\_w strengthens the suppression.

Ranking the shifts from the fiducial gives the following.

\begin{itemize}
  \item \textbf{At k=1:} the joint lowering of e\_w, v\_kin and eps\_kin (-2.4 pp) is the largest. It is followed by kappa\_w across its range (1.6 pp spread: +0.75 pp at kappa\_w=2, -0.86 pp at kappa\_w=4). Next is the joint raising of e\_w, v\_kin and eps\_kin (+0.33 pp), and last is M\_seed from 1.3 to 2.0 (+0.25 pp).
  \item \textbf{At k=3:} the joint lowering of e\_w, v\_kin and eps\_kin (about -8.8 pp) is again the largest shift sampled. It is followed by kappa\_w (5.98 pp spread; +2.9 pp at kappa\_w=2), then M\_seed (+1.0 pp), then the joint raising (+0.76 pp).
\end{itemize}
Of the parameters varied singly, kappa\_w has the largest one-at-a-time effect. However, it is not the largest shift sampled, because the joint move of the three AGN energy and kinetic parameters toward their lower bounds is larger.

At k=3 every shift exceeds the emulator 1$\sigma$ of 0.35--0.54\%. At k=1 the M\_seed shift and the joint-raising shift are only about 2$\sigma$.

e\_w, v\_kin and eps\_kin were only ever moved together, so their individual sensitivities are unknown. omega\_m and sigma\_8 were never varied. As a result, the roughly 1\% cosmology independence reported by FLAMINGO has not been checked for HACC.

For the kappa\_w and M\_seed runs, no values at k=0.1 or 0.5 h/Mpc are on the board. For the fiducial and strong points only, the large-scale suppression is about 0.03--0.04\% at k=0.1 h/Mpc and about 0.8--0.9\% at k=0.5 h/Mpc. The weak point is slightly enhanced at these k.

No comparison plot was produced for these runs.

\paragraph{Comparison with observational and calibrated benchmarks}
This is a range-of-model check, not a calibration. The cosmology and halo-mass weighting differ between HACC and the observational analyses, and HACC parameters do not map onto FLAMINGO fgas, FLAMINGO jet or BCemu parameters.

\textbf{At k=1 h/Mpc.} Across the six sampled points (fiducial, strong, weak, kappa\_w=2, kappa\_w=4, M\_seed=2), the suppression spans -0.2\% to 2.9\%. None reaches 5\% or 10\%. The most suppressed point, kappa\_w=2 at 2.92+/-0.15\%:

\begin{itemize}
  \item is about 0.5$\sigma$ from HSC-XXL;
  \item falls short of eROSITA+kSZ by about 3.5$\sigma$, using the observational error only;
  \item relative to Grandis et al., the fiducial and strong points sit about 1.2--1.4$\sigma$ below their 4.2\%.
\end{itemize}
\textbf{At k=3 h/Mpc.} The fiducial and strong HACC points give 8.0\% and 8.8\%. These are a factor of about 2.1--2.3 below the 18.4\% of arXiv:2309.02920, or about 3.1--3.4$\sigma$ against that paper's lower error bar. The largest sampled value, 10.9+/-0.5\% at kappa\_w=2, is still about 2.4$\sigma$ low.

This k=3 gap is unexplained, and the critic contested the comparison. Possible contributors are:

\begin{itemize}
  \item HACC's minimum lies beyond k=3.7 h/Mpc, while FLAMINGO's lies at 8--10 h/Mpc, so the curve shapes may differ.
  \item The cosmologies and halo-mass ranges differ. arXiv:2309.02920 is calibrated on haloes above about 3e13 Msun.
  \item HACC uses P\_grav as its reference spectrum.
  \item Combined extremes were not sampled, for example kappa\_w=2 together with high e\_w, v\_kin, eps\_kin and M\_seed.
\end{itemize}
\paragraph{Weak-lensing implications}
No C\_ell, xi\_+ or S8 shift was computed, because no tool provides them. All lensing statements in this report come from the literature.

The A\_mod-type suppression required by Amon \& Efstathiou and Preston et al. scales the whole non-linear part of P(k), with A\_mod$\approx$0.69 for KiDS. It reaches down to k$\approx$0.2 h/Mpc. The HACC fiducial and strong points suppress P(k) by less than 1\% at k$\leq$0.5 h/Mpc and by at most 2.5\% at k=1 h/Mpc.

Within the sampled range, the emulator is far from the suppression needed to reconcile shear with Planck. This comparison is approximate and was contested by the critic as a rough mapping between different parameterisations.

FLAMINGO likewise reports that the Amon--Efstathiou response lies outside the reach of the FLAMINGO emulator (arXiv:2410.17109). Both of its fgas-8$\sigma$ models, thermal and jet, fall inside the Bigwood et al. 2024 ranges, with or without kSZ.

The only sourced multipole range is that the baryon-sensitive shear extends to ell$\approx$5000 (arXiv:2410.12500). The DES Y3 numerical A\_mod was not read.

\subsubsection*{Discussion}
The emulator results agree qualitatively with the literature on two points:

\begin{itemize}
  \item Suppression grows with k through the readable range.
  \item Suppression is weak on linear scales, below 1\% at k\textless{}0.5 h/Mpc for the suppressed points. This is consistent with the \textless{}1\% bound at k\textless{}0.37 h/Mpc from arXiv:2309.02920.
\end{itemize}
Quantitatively, the sampled HACC models are mild-feedback models. At k=1 they sit at the low end of the HSC-XXL range. They cannot reproduce the strong feedback preferred by eROSITA+kSZ (arXiv:2512.02954) or the k=3 cluster-calibrated amplitude (arXiv:2309.02920).

The weak-AGN corner even shows a net enhancement of about 1\% near k$\approx$2.4 h/Mpc. The sensitivity scan suggests that the amplitude is controlled more by kappa\_w and by moving the energy and kinetic parameters toward their lower bounds than by raising them. Raising e\_w, v\_kin and eps\_kin together to their maximum adds only 0.3--0.8 pp. This saturation is reminiscent of FLAMINGO's reported saturation below fgas-8$\sigma$. However, the two parameterisations are not comparable, and this resemblance is not a demonstrated link.

There is an unresolved tension at k=3 that cannot be settled without the curve beyond k=3.7 h/Mpc. If HACC's minimum sits at k$\approx$8--10 h/Mpc like FLAMINGO's, the k=3 deficit reflects a genuinely weaker amplitude. If the shape differs, the comparison at a single k is misleading.

For weak lensing, the evidence assembled supports the FLAMINGO conclusion. Feedback consistent with X-ray gas fractions is not strong enough, nor large-scale enough, to remove the S8 tension. The KiDS and DES A\_mod solutions require suppression extending to k$\approx$0.2 h/Mpc, which is well beyond any of our sampled emulator points. The eROSITA+kSZ results favour stronger feedback than the fiducial simulations, which reopens part of this question. Even so, that paper's mitigation analysis treats baryons as a nuisance to calibrate, not as a resolution of the tension.

Dead ends:

\begin{itemize}
  \item \textbf{Reading the CSVs beyond k$\approx$3.73 h/Mpc.} The plain-text reader has no offset parameter and truncates each 799-bin CSV near k$\approx$3.73 h/Mpc. The table loader rejects the header row, and no file-writing tool was available to make a header-free copy. As a result, values at k=5 and 10 h/Mpc, the minimum and any upturn were never extracted.
  \item \textbf{Running other baryonic emulators.} No BCemu, BACCO or FLAMINGO emulator was available. Every computed number is from HACC.
\end{itemize}
\subsubsection*{Limitations}
\begin{itemize}
  \item \textbf{Truncated k range.} The k range effectively stops at 3.7 h/Mpc in our readout. The depth and position of the suppression minimum, which is a central quantity in the task, are not delivered. Whether the emulator itself reaches 10 h/Mpc is also not established.
  \item \textbf{Unverified mid- and high-k values.} The quoted evidence for every CSV ends near k$\approx$0.12 h/Mpc. Values at k$\geq$0.5 h/Mpc rest on run summaries.
  \item \textbf{Limited sampling.} Only six points were sampled, all at z=0 and at a single cosmology. There is no redshift scan, which matters because the lensing-weighted redshift is about 0.3--0.5. There is no cosmology variation, no fGas link and no combined-extreme points.
  \item \textbf{Missing literature numbers.} The FLAMINGO z=0 minimum depth and the DES Y3 A\_mod are unsourced. arXiv:2305.09827 and arXiv:2410.12500 were not added to the bibliography. The 1104.1174 claims about the minimum and upturn were never verified.
  \item \textbf{No lensing calculation.} No projected lensing calculation was possible.
  \item \textbf{Budget.} The two-wave budget was exhausted. The planned robustness and cosmology cross-checks were replaced by the limited sensitivity scan reported above.
\end{itemize}
\subsubsection*{Open directions}
\begin{enumerate}
  \item Re-read the 799-bin HACC Pk-ratio CSVs with a proper CSV reader to get the ratio at k=5 and 10 h/Mpc, the
\end{enumerate}
\subsection*{System note \small(2026-10-17 14:44:44)}
Report generated: /Users/nramachandra/hep-genesis-testing/projects/ramp-combo-baryonic-pk/reports/261017\_144243/report.pdf Title: Baryonic Suppression of the z=0 Matter Power Spectrum: HACC Emulator Predictions and Literature Assessment Supplement (implementation, tool-call log, board, extra figures, code index): /Users/nramachandra/hep-genesis-testing/projects/ramp-combo-baryonic-pk/reports/261017\_144243/supplement.pdf Transcript: /Users/nramachandra/hep-genesis-testing/projects/ramp-combo-baryonic-pk/reports/261017\_144243/transcript.pdf Folder: /Users/nramachandra/hep-genesis-testing/projects/ramp-combo-baryonic-pk/reports/261017\_144243 (report.tex, supplement.tex, transcript.tex, figures/, bundle.json)

\subsection*{User \small(2026-10-23 14:31:27)}
Survey recent work on baryonic suppression of the matter power spectrum. Use the hydrodynamics emulator tools to predict the Pk-ratio at z=0 for a couple of parameter points, explain the AGN feedback physics of suppression, and assess the impact on weak-lensing cosmology.

\subsection*{Assistant \small(research mode, 2026-10-23 14:31:27)}
\textit{Plan:}\begin{enumerate}
  \item (closed) Find the strongest suppression the HACC emulator can produce at z=0. Then measure how the Pk-ratio evolves from z=0 to z=1 for the fiducial point and the maximum-suppression point.
  \item (closed) Resolve the HACC Pk-ratio above k=3.7 h/Mpc (q-06, x-01), giving values at k=5 and 10 and the location and depth of the minimum. Also link the suppression to the group-scale gas fraction using the emulator's fGas observable.
  \item (closed) Close the remaining literature gaps: the DES Y3 A\_mod value (q-07), the depth of the FLAMINGO z=0 minimum (q-08), and cite entries for 2305.09827, 2410.12500, 2512.02954 and 2309.02920.
  \item (closed) Read the HACC outputs already written but never read (fGas for fiducial, strongAGN and weakAGN; maxsupp\_lowkin\_z0; maxsupp\_lowM\_z0; fiducial\_z1). Pair f\_gas near M500c \textasciitilde{} 1e14 Msun with S(k=1) and S(k=3) to test the van Daalen-type relation, and establish whether the design range can reach the eROSITA+kSZ level.
  \item (blocked) Verify the contested literature numbers verbatim and close the FLAMINGO gap. Quote 10+/-2\% and 5+/-4\% from 2512.02954, quote 0.042 and 0.184 at k=1 and 3 from 2309.02920, and quote the FLAMINGO z=0 S(k) at k=1 and the depth and location of its minimum from 2410.17109.
  \item (blocked) Make the weak-lensing impact quantitative. Extract one numerical S8 (or Omega\_m) bias from ignoring baryons, with its survey configuration. Compare HACC maxsupp S(k) at k=0.2-0.5 with the Preston et al. requirement that suppression extend to k\textasciitilde{}0.2 h/Mpc and with the size of A\_mod.
\end{enumerate}
\textit{Tool calls:} \texttt{list\_sterile\_nu\_skills}, \texttt{list\_observables}, \texttt{list\_agent\_files}, \texttt{list\_input\_files}, \texttt{list\_datasets}, \texttt{list\_files}, \texttt{get\_dispatch}, \texttt{search\_arxiv\_abstracts}, \texttt{search\_arxiv\_abstracts}, \texttt{search\_arxiv\_abstracts}, \texttt{search\_arxiv\_abstracts}, \texttt{search\_arxiv\_abstracts}, \texttt{search\_arxiv\_abstracts}, \texttt{search\_arxiv\_abstracts}, \texttt{search\_arxiv\_abstracts}, \texttt{search\_arxiv\_abstracts}, \texttt{search\_arxiv\_abstracts}, \texttt{load\_table\_data}, \texttt{describe\_parameters}, \texttt{list\_observables}, \texttt{list\_observables}, \texttt{list\_files}, \texttt{describe\_parameters}, \texttt{read\_text\_file}, \texttt{download\_arxiv\_full\_text}, \texttt{read\_arxiv\_chunk}, \texttt{get\_paper\_figures}, \texttt{predict\_observable}, \texttt{predict\_observable}, \texttt{predict\_observable}, \texttt{predict\_observable}, \texttt{predict\_observable}, \texttt{predict\_observable}, \texttt{download\_arxiv\_full\_text}, \texttt{download\_arxiv\_full\_text}, \texttt{get\_arxiv\_metadata}, \texttt{read\_arxiv\_chunk}, \texttt{read\_arxiv\_chunk}, \texttt{read\_arxiv\_chunk}, \texttt{read\_arxiv\_chunk}, \texttt{read\_arxiv\_chunk}, \texttt{read\_arxiv\_chunk}, \texttt{read\_arxiv\_chunk}, \texttt{post\_finding}, \texttt{read\_text\_file}, \texttt{plot\_observable\_comparison}, \texttt{plot\_observable\_comparison}, \texttt{post\_finding}, \texttt{post\_finding}, \texttt{post\_finding}, \texttt{post\_finding}, \texttt{post\_finding}, \texttt{post\_finding}, \texttt{post\_finding}, \texttt{post\_dead\_end}, \texttt{read\_arxiv\_chunk}, \texttt{read\_arxiv\_chunk}, \texttt{post\_finding}, \texttt{post\_dead\_end}, \texttt{read\_text\_file}, \texttt{predict\_observable}, \texttt{predict\_observable}, \texttt{predict\_observable}, \texttt{post\_finding}, \texttt{post\_finding}, \texttt{post\_dead\_end}, \texttt{post\_question}, \texttt{cite}, \texttt{cite}, \texttt{cite}, \texttt{cite}, \texttt{post\_finding}, \texttt{post\_finding}, \texttt{post\_question}, \texttt{download\_arxiv\_full\_text}, \texttt{read\_arxiv\_chunk}, \texttt{read\_arxiv\_chunk}, \texttt{read\_arxiv\_chunk}, \texttt{list\_files}, \texttt{list\_files}, \texttt{download\_arxiv\_full\_text}, \texttt{download\_arxiv\_full\_text}, \texttt{read\_arxiv\_chunk}, \texttt{read\_arxiv\_chunk}, \texttt{download\_arxiv\_full\_text}, \texttt{read\_arxiv\_chunk}, \texttt{read\_arxiv\_chunk}, \texttt{read\_arxiv\_chunk}, \texttt{post\_finding}, \texttt{post\_finding}, \texttt{post\_question}, \texttt{post\_finding}, \texttt{post\_finding}, \texttt{post\_question}, \texttt{post\_finding}, \texttt{post\_finding}, \texttt{post\_finding}, \texttt{post\_finding}, \texttt{post\_question}\par
\textbf{Survey} --- 52 tools on 5 servers, 7 read-only probes, 20 papers found.

Note first: I have no hydrodynamics or baryonic emulator tool. My only tools were arXiv and web search, so I did not compute any z=0 Pk-ratio (P\_hydro/P\_DMO) values. Any numbers for specific parameter points have to come from running BACCO (2011.15018), BCemu (2108.08863), HMcode-2020 (2009.01858) or the van Daalen et al. power-spectrum library (1906.00968) locally. I give no suppression amplitudes because I did not read results beyond titles and short abstracts. Physics: AGN feedback, and stellar feedback in lower-mass haloes, pushes gas out of group- and cluster-mass haloes (about 1e13-1e14.5 Msun). This lowers their baryon fraction and flattens their inner mass profiles, which suppresses P(k) at k of about 0.1-10 h/Mpc relative to dark-matter-only runs. At higher k, gas cooling and star formation push the ratio back up. Several groups (van Daalen 1906.00968, resummation model 2509.04552, Debackere 1908.05765, BAHAMAS/cosmo-OWLS 2309.06479, TNG 2511.10634) find the suppression is mainly set by the baryon fraction of group-scale haloes. That ties it to X-ray, tSZ and kSZ observables. For weak lensing, these scale-dependent changes bias S8 if left out (1105.1075, 1407.4301). Amon and Efstathiou (2206.11794, 2305.09827) argue that stronger-than-standard nonlinear suppression could reconcile KiDS and DES cosmic shear with Planck, and recent kSZ/X-ray results (FLAMINGO 2410.17109, 2602.02484) hint that feedback may be stronger than in fiducial simulations. Current mitigation options are scale cuts, marginalising over HMcode or baryonification parameters, theory-error covariances (2410.12500), and outside constraints from SZ, X-ray or FRBs (2301.02186, 2401.03510, 2501.17922).

Quantities worth checking: Pk-ratio S(k)=P\_hydro/P\_DMO at z=0 for k=0.1-20 h/Mpc at two or more feedback settings (e.g. HMcode log10 T\_AGN=7.6 vs 8.0; BACCO M\_c=1e13 vs 1e14.5 Msun/h); Scale and depth of the suppression minimum, and the k where S(k) first drops below 0.99; Halo baryon fraction at about 1e14 Msun vs the size of the suppression (van Daalen-type relation); Shift in S8 or Omega\_m when lensing data are fit with a DMO model vs the baryonic model; Size of suppression needed to reconcile lensing S8 with Planck (Amon and Efstathiou amplitude parameter A\_mod); Redshift dependence of S(k) between z=0 and z=1.

\textbf{Plan} --- 3 lines of investigation:

\begin{enumerate}
  \item Find the strongest suppression the HACC emulator can produce at z=0. Then measure how the Pk-ratio evolves from z=0 to z=1 for the fiducial point and the maximum-suppression point. --- Earlier runs varied the parameters one or a few at a time. The largest suppression seen at k=1 h/Mpc was 2.9\% (kappa\_w=2), and the combined extreme point was never sampled. The headline comparison is that HACC is a factor of about 2 below 2309.02920 at k=3 and cannot reach eROSITA+kSZ. That comparison is incomplete until we know the emulator's maximum. Redshift dependence is also on the survey's checklist and has not been done.
  \item Resolve the HACC Pk-ratio above k=3.7 h/Mpc (q-06, x-01), giving values at k=5 and 10 and the location and depth of the minimum. Also link the suppression to the group-scale gas fraction using the emulator's fGas observable. --- Every curve is still deepening at the point where the text reader truncates. That leaves the minimum, the high-k upturn and the comparison with the FLAMINGO minimum at k of about 8-10 h/Mpc unknown. The fGas link would give a van Daalen-type check of the physics, which is the AGN-mechanism part of the task.
  \item Close the remaining literature gaps: the DES Y3 A\_mod value (q-07), the depth of the FLAMINGO z=0 minimum (q-08), and cite entries for 2305.09827, 2410.12500, 2512.02954 and 2309.02920. --- The weak-lensing assessment and the comparison table depend on these. Several of them are currently stated without validated citations in references.bib.
Critic's view: revise --- Three corner-style points cannot establish the emulator's maximum suppression. Report the result as 'largest suppression among the sampled points'. Alternatively, spend the spare calls on a few extra corners, for example with kappa\_w and e\_w at their bounds and the kinetic parameters varied, before claiming that HACC cannot reach 10\% at k=1.; Step 1 must confirm, before any run, that every value lies inside the describe\_parameters bounds. This applies to e\_w=1.0, v\_kin=1.2, eps\_kin=1.2, M\_seed=2.0 and the cosmology omega\_m=0.14, sigma\_8=0.8. Points at the corners of the design hypercube carry the largest emulator error, so the 1-sigma must be quoted next to every ratio. Do not silently extrapolate.; Confirm the emulator's k units (h/Mpc or 1/Mpc) before comparing with HSC-XXL, eROSITA+kSZ and 2309.02920. Sigma comparisons should combine the emulator 1-sigma and the observational error in quadrature, and the method should be stated. The HSC-XXL and eROSITA+kSZ numbers (5+/-4\% and 10+/-2\%) currently have no source arXiv ID or validated citation. Add them to d-10, or drop the sigma claims..

\end{enumerate}
\textbf{Wave 1} --- now exploring, in parallel: Find the strongest suppression the HACC emulator can produce at z=0. Then measure how the Pk-ratio evolves from z=0 to z=1 for the fiducial point and the maximum-suppression point.; Resolve the HACC Pk-ratio above k=3.7 h/Mpc (q-06, x-01), giving values at k=5 and 10 and the location and depth of the minimum. Also link the suppression to the group-scale gas fraction using the emulator's fGas observable.; Close the remaining literature gaps: the DES Y3 A\_mod value (q-07), the depth of the FLAMINGO z=0 minimum (q-08), and cite entries for 2305.09827, 2410.12500, 2512.02954 and 2309.02920..

\textbf{Finding:} The HACC Pk-ratio emulator has trained redshifts z = 0, 0.1, 0.5, 1.0 and 2.0, with 799 k bins. Its design bounds are kappa\_w 2-4, e\_w 0.2-1.0, M\_seed 0.6-2.0, v\_kin 0.1-1.2, eps\_kin 0.02-1.2, omega\_m 0.12-0.155 and sigma\_8 0.7-0.9. Predictions were written for maxsupp, maxsupp\_lowkin and maxsupp\_lowM at z=0, and for fiducial and maxsupp at z=0 and at z=1.

\textbf{Finding:} HACC Pk-ratio (P\_hydro/P\_grav) at z=0 for the strongAGN point (kappa\_w=3.0, e\_w=1.0, M\_seed=1.3, v\_kin=1.2, eps\_kin=1.2, omega\_m=0.14, sigma\_8=0.8) is 0.851 at k=5.25, 0.8205 at k=7.01, 0.7946 at k=10.01 and 0.7893 at k=12.55 h/Mpc (the last tabulated k). The curve flattens to a plateau of about 0.789-0.790 for k of about 11.6-12.6 h/Mpc, so there is no resolved upturn: the minimum is at or just beyond k\_max of about 12.6 h/Mpc and has a depth of about 21\%.

\textbf{Finding:} The HACC Pk-ratio emulator k-grid (799 bins) runs from 0.022 to about 12.6 h/Mpc. That answers q-01: the grid reaches 10 h/Mpc but stops at about 12.6 h/Mpc, so the high-k upturn seen in FLAMINGO and OWLS-type runs is not covered. For the strongAGN point the ratio is still decreasing slowly (by less than 0.001 per 0.1 h/Mpc) at the edge.

\textbf{Finding:} At k of about 3.7 h/Mpc and z=0, the fiducial point (kappa\_w=3, e\_w=0.6, v\_kin=0.65, eps\_kin=0.61) has a Pk-ratio of 0.8984 (std 0.004), against 0.8909 for strongAGN. At k=1.01 the values are 0.9779 (fiducial) and 0.9746 (strongAGN), and at k=3.01 they are 0.9197 and 0.9121. Strong kinetic feedback therefore deepens the suppression by only about 0.3-0.8 percentage points over k=1-3.7 h/Mpc.

\textbf{Finding:} Siegel et al. (arXiv:2512.02954) find from joint eROSITA X-ray gas fractions plus SDSS/DESI+ACT kSZ profiles a matter power suppression of 10 +/- 2\% at k = 1 h/Mpc, stronger than the fiducial FLAMINGO and BAHAMAS models. HSC-XXL gas fractions prefer 5 +/- 4\% at k = 1 h/Mpc.

\textbf{Finding:} Grandis et al. (arXiv:2309.02920) use cluster gas and stellar fractions to infer a suppression of 0.042 (+0.012/-0.014) at k = 1 h/Mpc and 0.184 (+0.026/-0.031) at k = 3 h/Mpc. They predict, at 97.5\% credibility, that baryons change P(k) by less than 1\% for k \textless{} 0.37 h/Mpc, and on that basis question whether baryon feedback drives the cosmic-shear vs CMB S8 discrepancy.

\textbf{Finding:} Preston, Amon and Efstathiou (arXiv:2305.09827) model P\_m = P\_L + A\_mod (P\_NL - P\_L). Their KiDS-1000 Paper I value is A\_mod $\approx$ 0.69. For DES Y3 the abstract says only that the required suppression is ''not as strong as inferred using KiDS data, but is still more extreme than predictions from recent numerical simulations'', and that it must extend to k \textasciitilde{} 0.2 h/Mpc. The DES Y3 numerical A\_mod was not in the portion of the text read so far.

\textbf{Finding:} Maraio, Hall and Taylor (arXiv:2410.12500) find that for a Euclid-like survey with ell\_max = 5000, even multi-parameter baryon models (including HMCode-2020) can bias Omega\_m and S8 significantly. A theoretical-error covariance on P(k) reduces this bias to acceptable levels with only a modest increase in parameter variances. The abstract gives no numerical S8 shift.

\textbf{Finding:} At z=0 the combined extreme HACC point (kappa\_w=2, e\_w=1.0, M\_seed=2.0, v\_kin=1.2, eps\_kin=1.2, omega\_m=0.14, sigma\_8=0.8) gives P\_hydro/P\_grav = 0.99944+/-0.00003 at k=0.101, 0.9868+/-0.0008 at k=0.510, 0.9638+/-0.0020 at k=1.013, 0.9179+/-0.0045 at k=2.002, 0.8698+/-0.0065 at k=3.007 and 0.8400+/-0.0076 at k=3.698 h/Mpc. That is a suppression of 3.6+/-0.2\% at k=1 and 13.0+/-0.7\% at k=3. Against HSC-XXL (5+/-4\% at k=1) the gap is 0.35 sigma, so the point is consistent. Against eROSITA+kSZ (10+/-2\% at k=1) it falls short by 3.2 sigma. Against 2309.02920 (4.2\% at k=1, 18.4\% at k=3) it is 0.6 points low at k=1 (about 3 sigma of emulator error alone) and 5.4 points low at k=3 (about 8 sigma of emulator error alone). This point reaches neither 10\% at k=1 nor 18\% at k=3.

\textbf{Finding:} Answers q-07. Preston, Amon and Efstathiou (arXiv:2305.09827, Eq. 3) find that Planck LCDM is compatible with DES Y3 cosmic shear if non-linear power is suppressed with A\_mod = 0.858 +/- 0.052 (no scale cuts, Planck prior on S8 and Omega\_m). This is a weaker suppression than the KiDS-1000 value A\_mod $\approx$ 0.69 from Paper I (arXiv:2206.11794).

\textbf{Finding:} fGas (cluster gas fraction vs M500c, 3 mass bins) was predicted at z=0 for the fiducial, strongAGN and weakAGN points. The CSVs are written, but their values were not read because the tool budget ran out, so the van Daalen-type pairing of gas fraction with suppression is still open.

\textbf{Finding:} For the combined extreme HACC point, suppression weakens strongly from z=0 to z=1. At k=1.013 h/Mpc the ratio rises from 0.9638+/-0.0020 at z=0 to 0.9781+/-0.0018 at z=1, so the suppression drops from 3.6\% to 2.2\%. At k=3.007 it rises from 0.8698+/-0.0065 to 0.9299+/-0.0055 (13.0\% to 7.0\%). At k=3.70 it rises from 0.8400+/-0.0076 to 0.9214+/-0.0064 (16.0\% to 7.9\%). Within k of 3.7 h/Mpc the response therefore grows as z falls, which agrees qualitatively with FLAMINGO (2410.17109). The ratio is still falling at k=3.7 at both redshifts, so the minimum lies beyond the range read here and its shift with redshift could not be measured.

\textbf{Finding:} Four validated cite entries (arxiv230509827, arxiv241012500, arxiv251202954, arxiv230902920) were written to /Users/nramachandra/hep-genesis-testing/projects/ramp-combo-baryonic-pk/runs/261023\_143127/references.bib, each with a verbatim abstract quote.

\textbf{Completed --- Resolve the HACC Pk-ratio above k=3.7 h/Mpc (q-06, x-01), giving values at k=5 and 10 and the location and depth of the minimum. Also link the suppression to the group-scale gas fraction using the emulator's fGas observable..} The high-k part of this direction is resolved. The fGas link is only half done: the predictions were run but never read.

\begin{enumerate}
  \item Reading past k=3.7 h/Mpc. load\_table\_data cannot read the Pk-ratio CSVs because it does not skip the header row 'x,y\_mean,y\_std' and has no skiprows option (dead end x-04, tc-...-19). The files can be read in full with read\_arxiv\_chunk using start/length offsets.
\end{enumerate}
\begin{enumerate}
  \item Emulator k range (answers q-01, finding f-018). The HACC Pk-ratio grid has 799 bins from 0.022 to about 12.6 h/Mpc. It reaches 10 h/Mpc but nothing beyond about 12.6 h/Mpc, so any high-k upturn is outside the emulator range.
\end{enumerate}
\begin{enumerate}
  \item strongAGN point (kappa\_w=3, e\_w=1.0, v\_kin=1.2, eps\_kin=1.2, omega\_m=0.14, sigma\_8=0.8), z=0 (finding f-017). P\_hydro/P\_grav is 0.851 at k=5.25, 0.8205 at k=7.01, 0.7946 (std 0.0098) at k=10.01 and 0.7893 at k=12.55 h/Mpc. The curve flattens to a plateau of about 0.789-0.790 over k of about 11.6-12.6 h/Mpc. There is no upturn inside the range, so the minimum sits at or just past the grid edge with a depth of about 21\%. FLAMINGO's minimum is at k of about 8-10 h/Mpc; HACC's falls at k of 12 h/Mpc or higher, but it is not resolved.
\end{enumerate}
\begin{enumerate}
  \item Fiducial point vs strongAGN (finding f-019). Fiducial values are 0.9779 at k=1.01, 0.9197 at k=3.01 and 0.8984 at k=3.70 h/Mpc. The matching strongAGN values are 0.9746, 0.9121 and 0.8909, so strong kinetic feedback adds only about 0.3-0.8 percentage points of suppression over this range. The fiducial curve beyond k=3.7 and the kappaw2\_z0 and weakAGN\_z0 curves above k=3.7 were not extracted. Their files exist under runs/261017\_143153/ in folders of the same names, and their headers were checked.
\end{enumerate}
\begin{enumerate}
  \item Gas fraction link (finding f-026). fGas at z=0 (3 M500c bins) was predicted for the fiducial, strong and weak points. The output files are runs/261023\_143127/fgas\_fiducial/, fgas\_strong/ and fgas\_weak/, each containing fGas\_z0.0000\_prediction.csv. The values were not read before the tool budget ran out. That means the gas fraction at about 1e14 Msun is not yet paired with the k=1 and k=3 suppression, and whether HACC reproduces the van Daalen (1906.00968) relation is still unknown. This is posted as q-12.
\end{enumerate}
No figures were produced in this direction.

\textbf{Completed --- Find the strongest suppression the HACC emulator can produce at z=0. Then measure how the Pk-ratio evolves from z=0 to z=1 for the fiducial point and the maximum-suppression point..} The direction is mostly done, with three gaps. First, the fiducial z=0 and z=1 values were written to CSV but not read, so the z=0 vs z=1 comparison so far covers only one parameter point. Second, the maxsupp\_lowkin and maxsupp\_lowM variants were plotted but their values were not read, so it is still untested whether the kinetic parameters or M\_seed act non-monotonically. Third, only one of the two planned plots survives. Results: the Pk-ratio emulator is trained at z = 0, 0.1, 0.5, 1.0 and 2.0, and its design bounds are in f-016. At z=0 the combined extreme point (kappa\_w=2, e\_w=1.0, M\_seed=2.0, v\_kin=1.2, eps\_kin=1.2, omega\_m=0.14, sigma\_8=0.8) gives P\_hydro/P\_grav = 0.99944+/-0.00003 at k=0.10, 0.9868+/-0.0008 at k=0.51, 0.9638+/-0.0020 at k=1.01, 0.9179+/-0.0045 at k=2.0, 0.8698+/-0.0065 at k=3.0 and 0.8400+/-0.0076 at k=3.7 h/Mpc. That is a suppression of 3.6+/-0.2\% at k=1 and 13.0+/-0.7\% at k=3 (f-024). Comparisons: consistent with HSC-XXL (0.35 sigma); short of eROSITA+kSZ by 3.2 sigma; below 2309.02920 by 0.6 points at k=1 and 5.4 points at k=3, counting emulator error only. The extreme point reaches neither 10\% at k=1 nor 18\% at k=3. It is the strongest suppression found here, but the low-kinetic and low-M\_seed variants were not read out, so it is not shown to be the maximum over the whole design range. Redshift trend for the extreme point (f-027): suppression drops from z=0 to z=1, from 3.6\% to 2.2\% at k=1 and from 13.0\% to 7.0\% at k=3. This agrees with FLAMINGO's finding that the response grows as z falls. The ratio is still falling at k=3.7 at both redshifts, so the shift of the minimum could not be tested. Files: run dir /Users/nramachandra/hep-genesis-testing/projects/ramp-combo-baryonic-pk/runs/261023\_143127/, subdirectories maxsupp\_z0, maxsupp\_lowkin\_z0, maxsupp\_lowM\_z0, fiducial\_z0, fiducial\_z1 and maxsupp\_z1. Plot: Figure 28 = Pk-ratio\_comparison.png, showing z=0 vs z=1 for fiducial and maxsupp. The z=0 variant plot was overwritten because both plot calls write the same filename (dead end x-03). Fix: give each plot its own output\_dir subdirectory.

\textbf{Completed --- Close the remaining literature gaps: the DES Y3 A\_mod value (q-07), the depth of the FLAMINGO z=0 minimum (q-08), and cite entries for 2305.09827, 2410.12500, 2512.02954 and 2309.02920..} Two of the three targets are met. The FLAMINGO z=0 minimum depth (q-08) is still open.

q-07 answered (f-025): Preston, Amon and Efstathiou (arXiv:2305.09827, Eq. 3) give A\_mod = 0.858 +/- 0.052 for DES Y3 with no scale cuts and a Planck prior on S8 and Omega\_m. The verbatim quote is 'Amod = 0.858$\pm$ 0.052, DES Y3 (no scale cuts)' (tc-...-57). For comparison, KiDS-1000 in Paper I gave A\_mod $\approx$ 0.69 (tc-...-41). The same paper says the suppression must extend to k \textasciitilde{} 0.2 h/Mpc. Without A\_mod, DES Y3 S8 sits 0.9 sigma (variant 3) to 1.8 sigma (variant 4) below Planck S8 = 0.828 +/- 0.016.

Other results:

\begin{itemize}
  \item f-020: Siegel et al. (arXiv:2512.02954) find 10 +/- 2\% suppression at k = 1 h/Mpc from eROSITA plus kSZ data, stronger than fiducial FLAMINGO and BAHAMAS. HSC-XXL data give 5 +/- 4\%.
  \item f-021: Grandis et al. (arXiv:2309.02920) find 0.042 (+0.012/-0.014) at k = 1 h/Mpc and 0.184 (+0.026/-0.031) at k = 3 h/Mpc, and less than 1\% effect for k \textless{} 0.37 h/Mpc at 97.5\% credibility.
  \item f-023: Maraio, Hall and Taylor (arXiv:2410.12500) find that a theoretical-error covariance reduces the Omega\_m and S8 bias to acceptable levels for a Euclid-like survey at ell\_max = 5000. The part of the text read gives no numerical S8 shift.
  \item f-022: superseded by f-025.
\end{itemize}
Citations: four new validated entries were written to /Users/nramachandra/hep-genesis-testing/projects/ramp-combo-baryonic-pk/runs/261023\_143127/references.bib: arxiv230509827, arxiv241012500, arxiv251202954, arxiv230902920 (f-028).

Still open: q-08. get\_paper\_figures found no extractable figures for arXiv:2410.17109 (dead end x-02), so no FLAMINGO minimum depth is sourced to a figure. The follow-up is logged as q-13: read the FLAMINGO full text instead. The angular range of cosmic shear affected by baryons, from 2512.02954 and 2410.12500, was not reached within the budget.

\textbf{Critic, wave 1} --- continue: 8 findings stand up to the evidence, 5 questioned ('At z=0 the combined extreme HACC point (kappa\_w=2, e\_w=1.0, M\_seed=2.0, v\_kin=1.2, eps\_kin': The Pk-ratio values agree with the maxsupp CSV (tc-...-45). The Grandis et al. comparison uses only the emulator std (0.002 at k=1, 0.0065 at k=3) and leaves out the published errors (+0.012/-0.014 and +0.026/-0.031). Adding both in quadrature gives about 0.4-0.5 sigma at k=1 and about 1.7 sigma at k=3, not 3 and 8 sigma. The 3.2 sigma shortfall against eROSITA+kSZ rests on the 10+/-2\% value, which is not shown verbatim in the f-020 evidence.; 'At k of about 3.7 h/Mpc and z=0, the fiducial point (kappa\_w=3, e\_w=0.6, v\_kin=0.65, eps\_k': The strongAGN point changes e\_w (0.6 to 1.0) as well as v\_kin and eps\_kin. The 0.3-0.8 pp difference therefore cannot be credited to kinetic feedback alone. The quoted excerpts show only rows with k\textless{}0.12, so the k=1, 3 and 3.7 values are not visible in the evidence.; 'Siegel et al. (arXiv:2512.02954) find from joint eROSITA X-ray gas fractions plus SDSS/DES': The excerpt from tc-...-42 contains only the title, author list and page header. Neither 10+/-2\% (eROSITA+kSZ) nor 5+/-4\% (HSC-XXL) appears in it. f-024 depends on these numbers.; 'Grandis et al. (arXiv:2309.02920) use cluster gas and stellar fractions to infer a suppres': The excerpt from tc-...-36 ends partway through the abstract. The values 0.042, 0.184 and k\textless{}0.37 h/Mpc at 97.5\% credibility are not in the quoted text. These numbers are used quantitatively in f-024.).

Issues: The maximum suppression is not established. The maxsupp\_lowkin and maxsupp\_lowM CSVs exist but were never read. Cosmology was held at omega\_m=0.14 and sigma\_8=0.8, although the suppression depends on the cosmic baryon fraction. The kappa\_w effect is confounded with M\_seed: maxsupp (kappa\_w=2, M\_seed=2.0) and strongAGN (kappa\_w=3, M\_seed=1.3) differ in both. q-02 and q-10 remain unanswered.; The fGas CSVs were written but never read. The pairing of gas fraction with suppression is the central quantitative test of the AGN-physics explanation the task asks for, and reading the files costs about 3 tool calls.; The z=0 to z=1 trend has been measured for one point only; fiducial\_z1 was not read. Two FLAMINGO statements are unsourced while q-08/q-13 stays open: the 'agrees with FLAMINGO' claim, and the statement that FLAMINGO's minimum lies at k\textasciitilde{}8-10 h/Mpc..

\textbf{Next steps after wave 1} --- All earlier directions are closed and none is carried over. The final wave has about 11 tool calls, so it focuses on finishing work rather than starting new compute. First priority is reading the CSVs already on disk: fGas, lowkin, lowM and fiducial\_z1. That answers q-12, much of q-02 and q-10, and the critic's high-severity issues, and it tests the AGN-physics explanation through the f\_gas-to-suppression relation. Second is verbatim verification of the contested literature numbers (f-020, f-021), the FLAMINGO minimum (q-08, q-13), and a corrected error budget for f-024. Third, if budget remains, is one quantitative S8 bias from 2410.12500 and a test of HACC's large-scale suppression against the Preston et al. requirement, with the C\_ell forecast stated as future work. The critic's one-at-a-time parameter and cosmology scan and the P\_grav-definition check (q-09) are deferred because they cannot fit the budget; the report should list them as limitations. The task did not request facility runs, so the local and hosted execution used so far is sufficient.

New lines of investigation: Read the HACC outputs already written but never read (fGas for fiducial, strongAGN and weakAGN; maxsupp\_lowkin\_z0; maxsupp\_lowM\_z0; fiducial\_z1). Pair f\_gas near M500c \textasciitilde{} 1e14 Msun with S(k=1) and S(k=3) to test the van Daalen-type relation, and establish whether the design range can reach the eROSITA+kSZ level.; Verify the contested literature numbers verbatim and close the FLAMINGO gap. Quote 10+/-2\% and 5+/-4\% from 2512.02954, quote 0.042 and 0.184 at k=1 and 3 from 2309.02920, and quote the FLAMINGO z=0 S(k) at k=1 and the depth and location of its minimum from 2410.17109.; Make the weak-lensing impact quantitative. Extract one numerical S8 (or Omega\_m) bias from ignoring baryons, with its survey configuration. Compare HACC maxsupp S(k) at k=0.2-0.5 with the Preston et al. requirement that suppression extend to k\textasciitilde{}0.2 h/Mpc and with the size of A\_mod..

\textbf{Wave 2} --- now exploring, in parallel: Read the HACC outputs already written but never read (fGas for fiducial, strongAGN and weakAGN; maxsupp\_lowkin\_z0; maxsupp\_lowM\_z0; fiducial\_z1). Pair f\_gas near M500c \textasciitilde{} 1e14 Msun with S(k=1) and S(k=3) to test the van Daalen-type relation, and establish whether the design range can reach the eROSITA+kSZ level.; Verify the contested literature numbers verbatim and close the FLAMINGO gap. Quote 10+/-2\% and 5+/-4\% from 2512.02954, quote 0.042 and 0.184 at k=1 and 3 from 2309.02920, and quote the FLAMINGO z=0 S(k) at k=1 and the depth and location of its minimum from 2410.17109.; Make the weak-lensing impact quantitative. Extract one numerical S8 (or Omega\_m) bias from ignoring baryons, with its survey configuration. Compare HACC maxsupp S(k) at k=0.2-0.5 with the Preston et al. requirement that suppression extend to k\textasciitilde{}0.2 h/Mpc and with the size of A\_mod..

\textbf{Finding:} Corrected deficit significance for the HACC extreme point (f-024), with emulator and published errors combined in quadrature and the published error taken on the side facing the HACC value. Against Grandis et al. (2309.02920): at k=1, (4.2-3.6)/sqrt(0.2\textasciicircum{}2+1.4\textasciicircum{}2) = 0.4 sigma; at k=3, (18.4-13.0)/sqrt(0.7\textasciicircum{}2+3.1\textasciicircum{}2) = 1.7 sigma. Against eROSITA+kSZ (2512.02954): (10-3.6)/sqrt(0.2\textasciicircum{}2+2\textasciicircum{}2) = 3.2 sigma. Against HSC-XXL: 1.4/sqrt(0.2\textasciicircum{}2+4\textasciicircum{}2) = 0.35 sigma. This replaces the f-024 statements of about 3 sigma and 8 sigma, which counted emulator error only.

\textbf{Finding:} Full texts of FLAMINGO (2410.17109), Siegel et al. (2512.02954) and Grandis et al. (2309.02920) were downloaded to /Users/nramachandra/hep-genesis-testing/projects/ramp-combo-baryonic-pk/runs/261023\_143127/2410.17109.txt, 2512.02954.txt and 2309.02920.txt. They have not been read yet, so no verbatim quote has been checked and no FLAMINGO S(k) value is sourced.

\textbf{Finding:} Maraio, Hall and Taylor (arXiv:2410.12500, Table 3) work with an idealised Euclid-like (Stage-IV) cosmic-shear survey at ell\_max = 5000 and use Illustris as the ground truth. When the one-parameter HMCode baryon model is fitted without a theory error, Omega\_m is biased by 100*Delta Omega\_m = -1.317, which is -5.690 sigma. Adding the theoretical-error covariance cuts this to 0.128 (0.249 sigma). With Bahamas as the ground truth the bias is only 0.057 (0.258 sigma). Eagle with the 6-parameter model is still biased with the theory error included: 0.573 (1.459 sigma). In the text, the one-parameter model without theory error gives 'at least a 3 sigma offset in either Omega\_m or S8 for all hydro-sims excluding Bahamas'. For Simba, even ell\_max = 1000 is biased by more than 2 sigma in all cosmological parameters.

\textbf{Finding:} For the combined extreme point, the HACC maxsupp z=0 Pk-ratio keeps falling past k=3.7 h/Mpc: 0.8284+/-0.0079 at k=4.00, 0.7942+/-0.0089 at k=5.00, 0.7680+/-0.0095 at k=6.0, 0.7479+/-0.0100 at k=7.01 and 0.7356 at k=7.84 h/Mpc. That is already deeper than the strongAGN plateau of about 0.789 at k of about 12 h/Mpc (f-017), so kappa\_w=2 with M\_seed=2.0 gives more small-scale suppression than kappa\_w=3. The values at k=0.2-0.3 h/Mpc were not read within the budget.

\textbf{Blocked --- Verify the contested literature numbers verbatim and close the FLAMINGO gap. Quote 10+/-2\% and 5+/-4\% from 2512.02954, quote 0.042 and 0.184 at k=1 and 3 from 2309.02920, and quote the FLAMINGO z=0 S(k) at k=1 and the depth and location of its minimum from 2410.17109..} (The tool budget ran out right after the three full-text downloads, before any read\_arxiv\_chunk call. No verbatim quotes or FLAMINGO S(k) values were extracted. Next step: read the downloaded .txt files (q-14).) Only the error-budget part of this direction is done. No text was read, so f-020 and f-021 are still not confirmed verbatim and no FLAMINGO S(k) value has been quoted. (1) Full texts of 2410.17109 (FLAMINGO), 2512.02954 and 2309.02920 were downloaded to /Users/nramachandra/hep-genesis-testing/projects/ramp-combo-baryonic-pk/runs/261023\_143127/ as 2410.17109.txt, 2512.02954.txt and 2309.02920.txt (f-030). (2) The deficit significance for the HACC extreme point was recomputed (f-029). Emulator and published errors were combined in quadrature, using the published error on the side facing the HACC value. Against Grandis et al. the deficit is 0.4 sigma at k=1 h/Mpc (4.2\% vs 3.6\%) and 1.7 sigma at k=3 h/Mpc (18.4\% vs 13.0\%). Against eROSITA+kSZ it is 3.2 sigma at k=1, and against HSC-XXL it is 0.35 sigma. This replaces the f-024 values of about 3 and 8 sigma, which used emulator error only. The k=1 result is 0.4 sigma, close to the 0.5 sigma the lead expected. These sigmas still rest on the board values in f-020 and f-021, which have not been checked against the source text. The full-text reading is logged as q-14. f-022 was not formally retired, but f-025 already supersedes it.

\textbf{Blocked --- Make the weak-lensing impact quantitative. Extract one numerical S8 (or Omega\_m) bias from ignoring baryons, with its survey configuration. Compare HACC maxsupp S(k) at k=0.2-0.5 with the Preston et al. requirement that suppression extend to k\textasciitilde{}0.2 h/Mpc and with the size of A\_mod..} (The tool budget ran out before I could read the HACC values at k=0.2-0.3 h/Mpc and the Delta S8 table (Table 4) of 2410.12500.) Only part of this direction got done. I have one numerical cosmology bias from ignoring baryons, but it is an Omega\_m bias, not an S8 bias, and I did not read the HACC values at k=0.2-0.3 h/Mpc.

What was established:

\begin{enumerate}
  \item Omega\_m bias from Maraio, Hall and Taylor (arXiv:2410.12500, Table 3), read verbatim (f-031; tc-...-86). Survey configuration: an idealised Euclid-like Stage-IV cosmic-shear survey with ell\_max = 5000.
\end{enumerate}
\begin{itemize}
  \item Illustris as ground truth, analysed with the one-parameter HMCode model and no theory error: 100*Delta Omega\_m = -1.317, which is -5.69 sigma.
  \item The same case with the theoretical-error covariance added: 0.128, or 0.249 sigma.
  \item Bahamas as ground truth with the one-parameter model: 0.057, or 0.258 sigma.
  \item Eagle with the six-parameter model and the theory error included: still 0.573, or 1.459 sigma.
  \item The text says the one-parameter model without theory error gives at least a 3 sigma offset in either Omega\_m or S8 for every hydro simulation except Bahamas. For Simba as ground truth, even ell\_max = 1000 is biased by more than 2 sigma in all parameters.
  \item The full text is saved at runs/261023\_143127/2410.12500.txt. The Delta S8 values are in Table 4, past character 55000 of that file, and I did not read them.
\end{itemize}
\begin{enumerate}
  \item HACC maxsupp z=0 Pk-ratio beyond k=3.7 h/Mpc (f-032; tc-...-83). Values are 0.8284 at k=4.0, 0.7942 at k=5.0, 0.7680 at k=6.0, 0.7479 at k=7.0 and 0.7356 at k=7.84 h/Mpc. This is already deeper than the strongAGN plateau of about 0.789 (f-017).
\end{enumerate}
\begin{itemize}
  \item Source file: runs/261023\_143127/maxsupp\_z0/Pk-ratio\_z0.0000\_prediction.csv.
\end{itemize}
\begin{enumerate}
  \item HACC suppression versus the A\_mod requirement: this comparison is qualitative only, since the k=0.2-0.3 values are unread.
\end{enumerate}
\begin{itemize}
  \item From f-024, S(k) goes from 0.9994 at k=0.10 to 0.987 at k=0.51 h/Mpc, so at k\textasciitilde{}0.2 it is presumably below 1\%.
  \item That is consistent with Grandis et al.'s bound of less than 1\% for k \textless{} 0.37 h/Mpc.
  \item It is far weaker than the suppression that A\_mod = 0.858+/-0.052 (DES Y3) or about 0.69 (KiDS-1000) would need at k\textasciitilde{}0.2 h/Mpc.
  \item On this basis, HACC-level feedback is unlikely to resolve S8. That conclusion rests on interpolation, not on read-out values.
\end{itemize}
No figures were produced.

Open (q-15): read the maxsupp CSV for k=0.2-0.3 h/Mpc (roughly characters 0-3000), read Table 4 (Delta S8) of 2410.12500, and run a Limber C\_ell / S8 Fisher forecast using HACC S(k), which remains future work (q-03, q-11).

\textbf{Finding:} The HACC fGas emulator at z=0 has 3 M500c bins: 5.01e13, 7.94e13 and 1.26e14 Msun/h. In the bin nearest 1e14 (1.26e14 Msun/h), f\_gas = 0.1214+/-0.0010 for fiducial, 0.1221+/-0.0019 for strongAGN and 0.1414+/-0.0023 for weakAGN (e\_w=0.2, v\_kin=0.1, eps\_kin=0.02). At 7.94e13 Msun/h the values are 0.1090, 0.1098 and 0.1375.

\textbf{Finding:} Partial answer to q-12. Only two fGas points have S(k) read so far. Fiducial: f\_gas(1.26e14) = 0.1214+/-0.0010, with S(k=1) = 0.9779 and S(k=3) = 0.9197. StrongAGN: f\_gas = 0.1221+/-0.0019, with S(k=1) = 0.9746 and S(k=3) = 0.9121. The two f\_gas values agree within errors; strongAGN is nominally higher, yet its suppression is 0.3-0.8 points deeper. Between these two points there is therefore no resolvable van Daalen-type relation, and the nominal trend runs the wrong way. WeakAGN has a much higher f\_gas (0.1414), but its S(k) was not read in this direction, so a three-point test is still open. The strongAGN vs fiducial comparison is confounded: e\_w changes (0.6 to 1.0) at the same time as v\_kin and eps\_kin.

\textbf{Finding:} The HACC maxsupp\_lowkin point at z=0 (kappa\_w=2, e\_w=1.0, M\_seed=2.0, v\_kin=0.1, eps\_kin=0.02, omega\_m=0.14, sigma\_8=0.8) shows essentially no suppression. P\_hydro/P\_grav is 1.0002 at k=0.212, 1.0004+/-0.0007 at k=0.510, 1.0006+/-0.0019 at k=1.013, 1.0033+/-0.0042 at k=2.002, 0.9972+/-0.0062 at k=3.007 and 0.9883+/-0.0071 at k=3.698 h/Mpc. So in this emulator the kinetic AGN parameters (v\_kin, eps\_kin) drive almost all of the suppression at k up to 3.7 h/Mpc, even with e\_w at its maximum.

\textbf{Finding:} The HACC maxsupp\_lowM point at z=0 (kappa\_w=2, e\_w=1.0, M\_seed=0.6, v\_kin=1.2, eps\_kin=1.2, omega\_m=0.14, sigma\_8=0.8) suppresses more than the combined maxsupp point (M\_seed=2.0). P\_hydro/P\_grav is 0.9967 at k=0.212, 0.9835+/-0.0011 at k=0.510, 0.9557+/-0.0029 at k=1.013, 0.9026+/-0.0064 at k=2.002, 0.8521+/-0.0094 at k=3.007 and 0.8228+/-0.0109 at k=3.698 h/Mpc. The maxsupp values are 0.9638, 0.8698 and 0.8400 at k=1, 3 and 3.7. The suppression is 4.4+/-0.3\% at k=1 and 14.8+/-0.9\% at k=3, so maxsupp is not the strongest point, and low M\_seed increases the suppression. This is still short of eROSITA+kSZ (10+/-2\% at k=1, arXiv:2512.02954) by about 2.8 sigma. It is consistent with HSC-XXL (5+/-4\%) and with Grandis et al. at k=1 (4.2\%). At k=3 it is 3.6 points below Grandis (18.4\%). Among the points read, the HACC design range does not reach the eROSITA+kSZ level at k=1.

\textbf{Completed --- Read the HACC outputs already written but never read (fGas for fiducial, strongAGN and weakAGN; maxsupp\_lowkin\_z0; maxsupp\_lowM\_z0; fiducial\_z1). Pair f\_gas near M500c \textasciitilde{} 1e14 Msun with S(k=1) and S(k=3) to test the van Daalen-type relation, and establish whether the design range can reach the eROSITA+kSZ level..} Direction partly met. The f\_gas values and the two maxsupp variants were read; the weakAGN pairing and the fiducial z=0 vs z=1 comparison were not. (1) fGas at z=0 (f-033), in the bin nearest 1e14 Msun/h (M500c = 1.26e14 Msun/h): fiducial 0.1214+/-0.0010, strongAGN 0.1221+/-0.0019, weakAGN 0.1414+/-0.0023. (2) Gas-fraction vs suppression (f-034): this pairing is not met. Only 2 points have S(k) values (fiducial: S(k=1)=0.9779, S(k=3)=0.9197; strongAGN: 0.9746, 0.9121). Their f\_gas values agree within errors, and strongAGN is nominally higher despite its deeper suppression, so this pair shows no van Daalen-type relation. The comparison is also confounded, because e\_w changes together with the kinetic parameters. The weakAGN S(k) was not read, so the 3-point test is still open. (3) maxsupp\_lowkin (v\_kin=0.1, eps\_kin=0.02) shows essentially no suppression: S = 1.0006+/-0.0019 at k=1.013, 0.9972+/-0.0062 at k=3.007 and 0.9883 at k=3.70 h/Mpc. The kinetic parameters therefore drive the suppression (f-035). (4) maxsupp\_lowM (M\_seed=0.6) is stronger than maxsupp: S = 0.9967 at k=0.21, 0.9835 at k=0.51, 0.9557+/-0.0029 at k=1.013, 0.9026 at k=2.0, 0.8521+/-0.0094 at k=3.007 and 0.8228+/-0.0109 at k=3.70 h/Mpc. That is 4.4\% suppression at k=1 and 14.8\% at k=3 (f-036). (5) Of the points read, none reaches the eROSITA+kSZ level of 10+/-2\% at k=1. The strongest, lowM, falls short by about 2.8 sigma; it is consistent with HSC-XXL and with Grandis et al. at k=1, but 3.6 points low at k=3. (6) Not done (q-16): fiducial\_z1 at k=1 and k=3, weakAGN S(k), and a verbatim quote of the k\textasciitilde{}12.5 rows for f-018. Files: runs/261023\_143127/fgas\_\{fiducial,strong,weak\}/fGas\_z0.0000\_prediction.csv, maxsupp\_lowkin\_z0/ and maxsupp\_lowM\_z0/Pk-ratio\_z0.0000\_prediction.csv. No figures were produced.

\textbf{Critic, wave 2} --- continue: 5 findings stand up to the evidence, 3 questioned ('Corrected deficit significance for the HACC extreme point (f-024), with emulator and publi': The arithmetic checks out: 0.6/1.414 = 0.42, 5.4/3.18 = 1.70, 6.4/2.01 = 3.18 and 1.4/4.005 = 0.35. However, every input comes from unverified values. The 2512.02954 and 2309.02920 excerpts are cut off in the abstract or title block and do not contain 10+/-2\%, 5+/-4\%, 0.042 or 0.184. f-020 and f-021 are themselves contested. The 3.6\% emulator value at k=1 is also not in the quoted maxsupp excerpt, which stops at k=0.117. The sigma values cannot be treated as established until q-14 is done.; 'Maraio, Hall and Taylor (arXiv:2410.12500, Table 3) work with an idealised Euclid-like (St': The cited excerpt (tc-...-86, chars 40000-55000) shows only text about the covariance discussion. It does not contain Table 3 or any of the claimed numbers (-1.317, -5.690 sigma, 0.128, 0.057, 0.573, 1.459 sigma), nor the 'at least a 3 sigma offset' quote. The claim is marked high confidence and says the values were read verbatim, but the evidence as posted does not show them. A verbatim excerpt of the table rows is needed.; 'Partial answer to q-12. Only two fGas points have S(k) read so far. Fiducial: f\_gas(1.26e1': The f\_gas values are supported by the excerpts. The paired S(k=1) and S(k=3) values (0.9779/0.9197 and 0.9746/0.9121) are not in the quoted Pk excerpts, which stop at k of about 0.1, and they reuse f-019, which is contested on the board. The conclusion also rests on two points that agree within errors and differ in e\_w, v\_kin and eps\_kin together. That cannot support any statement about a van Daalen-type relation, including that the trend runs 'the wrong way'.).

Issues: Every literature-comparison significance on the board (f-024, f-029, f-036) depends on f-020 and f-021, and neither has been checked verbatim. The full texts are already on disk, so checking them is cheap. Until it is done, the headline conclusion that HACC falls short of eROSITA+kSZ by about 3 sigma is provisional.; The emulator values for lowM are supported. The comparisons to Siegel et al. and Grandis et al. inherit the contested status of f-020 and f-021. The claim that 'the design range does not reach eROSITA+kSZ' has only been tested at the kappa\_w=2, e\_w=1 corners with omega\_m=0.14 and sigma\_8=0.8. The lowM curve is still falling at k=3.7, so its minimum is unresolved. The cosmology axes (omega\_m 0.12-0.155, sigma\_8 0.7-0.9) were not explored, and they can shift the response through the baryon fraction and halo abundance.; There is a physically interesting signal here that the finding does not develop. Strong kinetic feedback leaves f\_gas(\textless{}R500c) at 1.26e14 essentially unchanged (0.1214 vs 0.1221), yet it deepens S(k). That suggests the HACC suppression is driven by gas removal beyond R500c, or by other halo masses, rather than by the R500c gas fraction. A controlled single-parameter sweep is needed to test this..

Stopping after wave 2: budget reached (2/2 waves used). Writing the report from the board.

\textit{Research report:}\par
\subsubsection*{Baryonic suppression of the matter power spectrum: HACC emulator predictions at z = 0, AGN-feedback physics and weak-lensing impact}
\subsubsection*{Abstract}
This report combines a literature survey of baryonic suppression of the matter power spectrum with predictions from a HACC-based hydrodynamics emulator. The emulator returns the ratio S(k) = P\_hydro/P\_grav, together with cluster gas fractions, over a 7-parameter design space: five AGN and black-hole parameters plus omega\_m and sigma\_8.

Main results from the emulator:

\begin{itemize}
  \item At z = 0 the strongest suppression among the points sampled is 4.4 $\pm$ 0.3\% at k = 1 h/Mpc and 14.8 $\pm$ 0.9\% at k = 3 h/Mpc. It comes from maximal kinetic feedback, maximal e\_w, kappa\_w = 2 and the lowest seed mass.
  \item Kinetic AGN feedback produces almost all of the suppression. With the kinetic parameters set low, S(k) stays within about 1\% of unity up to k = 3.7 h/Mpc.
  \item The suppression weakens substantially between z = 0 and z = 1. For the one point measured, it drops from 13.0\% to 7.0\% at k = 3 h/Mpc.
  \item The emulator's k-grid ends at about 12.6 h/Mpc, so the suppression minimum and any high-k upturn are not resolved.
  \item None of the points sampled reaches the 10 $\pm$ 2\% suppression at k = 1 h/Mpc reported from eROSITA plus kSZ data. The shortfall is about 2.8 to 3.2 sigma, but this rests on literature values that have not been verified verbatim.
  \item Two points with very different S(k) have the same gas fraction at M500c $\approx$ 1.26e14 Msun/h, so the van Daalen-type relation between gas fraction and suppression is not reproduced by the pair tested.
\end{itemize}
For weak lensing, the evidence is taken from the literature rather than our own forecast:

\begin{itemize}
  \item An Omega\_m bias of -5.7 sigma arises for a Euclid-like survey when Illustris-like feedback is fitted with a one-parameter HMCode model. A theoretical-error covariance removes most of this bias.
  \item Reconciling DES Y3 with Planck needs A\_mod = 0.858 $\pm$ 0.052. The suppression this implies at k $\approx$ 0.2 h/Mpc appears much larger than anything the emulator produces there.
\end{itemize}
\subsubsection*{Background}
\textbf{Physics of the suppression.} AGN feedback, with stellar feedback in lower-mass haloes, drives gas out of group- and cluster-mass haloes (about 1e13 to 1e14.5 Msun). This lowers their baryon fractions and flattens their inner mass profiles. The result is a suppression of P(k) relative to dark-matter-only runs at k $\approx$ 0.1--10 h/Mpc. At higher k, cooling and star formation push the ratio back up.

\textbf{Link to halo baryon fractions.} Several studies find that the suppression is set mainly by the baryon fraction of group-scale haloes:

\begin{itemize}
  \item the van Daalen et al. library of 92 simulation power-spectrum ratios (arXiv:1906.00968);
  \item halo-model work tying observed gas profiles to P(k) (arXiv:1908.05765);
  \item a BAHAMAS/cosmo-OWLS breakdown by halo mass (arXiv:2309.06479);
  \item an IllustrisTNG study of which halo masses and radii matter at k \textgreater{} 1 h/Mpc (arXiv:2511.10634);
  \item the FLAMINGO resummation model, which predicts suppression from observed halo baryon fractions (arXiv:2509.04552).
\end{itemize}
FLAMINGO also studies the baryon response across feedback variants calibrated to cluster gas fractions (arXiv:2410.17109). This link makes X-ray, tSZ and kSZ observations direct probes of the suppression.

\textbf{Public modelling tools.} These include the BACCO baryonification emulator (arXiv:2011.15018), BCemu (arXiv:2108.08863), HMcode-2020 with its log10 T\_AGN parameter (arXiv:2009.01858), an analytic redshift-independent formulation (arXiv:2504.15633), and the CAMELS simulation suite (arXiv:2010.00619).

\textbf{Observational constraints.}

\begin{itemize}
  \item Siegel et al. (arXiv:2512.02954) combine eROSITA gas fractions with kSZ profiles and infer 10 $\pm$ 2\% suppression at k = 1 h/Mpc. This is stronger than fiducial FLAMINGO and BAHAMAS. HSC-XXL gas fractions give 5 $\pm$ 4\%.
  \item Grandis et al. (arXiv:2309.02920) infer 0.042 (+0.012/-0.014) at k = 1 h/Mpc and 0.184 (+0.026/-0.031) at k = 3 h/Mpc from cluster gas and stellar fractions. They predict a change of less than 1\% for k \textless{} 0.37 h/Mpc at 97.5\% credibility.
\end{itemize}
Neither set of numbers has yet been checked against the full text; both come from abstract-level reading.

\textbf{Weak lensing.} Ignoring these scale-dependent changes biases S8 (arXiv:1105.1075). Feedback is also degenerate with neutrino mass (arXiv:1407.4301). Amon and Efstathiou (arXiv:2206.11794) and Preston, Amon and Efstathiou (arXiv:2305.09827) model P\_m = P\_L + A\_mod (P\_NL - P\_L). They argue that stronger-than-standard nonlinear suppression can reconcile cosmic shear with Planck. Further evidence for strong feedback is cited in arXiv:2602.02484.

Mitigation strategies in the literature are:

\begin{itemize}
  \item scale cuts;
  \item marginalising over baryonic parameters;
  \item theoretical-error covariances (arXiv:2410.12500);
  \item external SZ, kSZ or FRB constraints (arXiv:2301.02186, arXiv:2401.03510, arXiv:2501.17922).
\end{itemize}
\subsubsection*{Results}
\paragraph{Emulator configuration}
The Pk-ratio emulator has these properties:

\begin{itemize}
  \item \textbf{Redshifts:} trained at z = 0, 0.1, 0.5, 1.0 and 2.0.
  \item \textbf{k-grid:} 799 bins from 0.022 to about 12.6 h/Mpc. The k units are h/Mpc according to the file headers.
  \item \textbf{Parameter bounds:}
\end{itemize}
\begin{itemize}
  \item \textbf{Reference spectrum:} the ratio is taken relative to P\_grav, a gravity-only run. Whether this is equivalent to the P\_DMO used in the literature has not been checked.
\end{itemize}
\begin{longtable}{@{}p{0.480\textwidth}p{0.480\textwidth}@{}}
\toprule
\textbf{Parameter} & \textbf{Range} \\ \midrule \endhead
kappa\_w & 2--4 \\
e\_w & 0.2--1.0 \\
M\_seed & 0.6--2.0 (units of 1e6 Msun) \\
v\_kin & 0.1--1.2 \\
eps\_kin & 0.02--1.2 \\
omega\_m & 0.12--0.155 \\
sigma\_8 & 0.7--0.9 \\
\bottomrule\end{longtable}
All parameter points below lie within these bounds. Their full parameter labels, read from the prediction-file headers, are:

\begin{longtable}{@{}p{0.160\textwidth}p{0.160\textwidth}p{0.160\textwidth}p{0.160\textwidth}p{0.160\textwidth}p{0.160\textwidth}@{}}
\toprule
\textbf{Point} & \textbf{kappa\_w} & \textbf{e\_w} & \textbf{M\_seed} & \textbf{v\_kin} & \textbf{eps\_kin} \\ \midrule \endhead
fiducial & 3 & 0.6 & 1.3 & 0.65 & 0.61 \\
strongAGN & 3 & 1.0 & 1.3 & 1.2 & 1.2 \\
weakAGN & 3 & 0.2 & 1.3 & 0.1 & 0.02 \\
maxsupp & 2 & 1.0 & 2.0 & 1.2 & 1.2 \\
maxsupp\_lowkin & 2 & 1.0 & 2.0 & 0.1 & 0.02 \\
maxsupp\_lowM & 2 & 1.0 & 0.6 & 1.2 & 1.2 \\
\bottomrule\end{longtable}
Every file header lists omega\_m = 0.14 and sigma\_8 = 0.8, so the cosmology axes were never varied.

\paragraph{Pk-ratio at z = 0 for the sampled feedback points}
\textbf{maxsupp (combined extreme corner).} S(k) at z = 0:

\begin{longtable}{@{}p{0.480\textwidth}p{0.480\textwidth}@{}}
\toprule
\textbf{k (h/Mpc)} & \textbf{S(k)} \\ \midrule \endhead
0.101 & 0.99944 $\pm$ 0.00003 \\
0.510 & 0.9868 $\pm$ 0.0008 \\
1.013 & 0.9638 $\pm$ 0.0020 \\
2.002 & 0.9179 $\pm$ 0.0045 \\
3.007 & 0.8698 $\pm$ 0.0065 \\
3.698 & 0.8400 $\pm$ 0.0076 \\
4.00 & 0.8284 $\pm$ 0.0079 \\
5.00 & 0.7942 $\pm$ 0.0089 \\
6.0 & 0.7680 $\pm$ 0.0095 \\
7.01 & 0.7479 $\pm$ 0.0100 \\
7.84 & 0.7356 \\
\bottomrule\end{longtable}
This is 3.6 $\pm$ 0.2\% suppression at k = 1 and 13.0 $\pm$ 0.7\% at k = 3 h/Mpc.

\textbf{maxsupp\_lowM (M\_seed = 0.6).} Lowering the seed mass from 2.0 to 0.6 deepens the suppression:

\begin{longtable}{@{}p{0.480\textwidth}p{0.480\textwidth}@{}}
\toprule
\textbf{k (h/Mpc)} & \textbf{S(k)} \\ \midrule \endhead
0.212 & 0.9967 \\
0.510 & 0.9835 $\pm$ 0.0011 \\
1.013 & 0.9557 $\pm$ 0.0029 \\
2.002 & 0.9026 $\pm$ 0.0064 \\
3.007 & 0.8521 $\pm$ 0.0094 \\
3.698 & 0.8228 $\pm$ 0.0109 \\
\bottomrule\end{longtable}
This is 4.4 $\pm$ 0.3\% at k = 1 and 14.8 $\pm$ 0.9\% at k = 3 h/Mpc, the largest suppression among the sampled points. It is not shown to be the maximum over the design range.

\textbf{maxsupp\_lowkin (kinetic parameters at minimum).} This point shows essentially no suppression:

\begin{longtable}{@{}p{0.480\textwidth}p{0.480\textwidth}@{}}
\toprule
\textbf{k (h/Mpc)} & \textbf{S(k)} \\ \midrule \endhead
0.212 & 1.0002 \\
0.510 & 1.0004 $\pm$ 0.0007 \\
1.013 & 1.0006 $\pm$ 0.0019 \\
2.002 & 1.0033 $\pm$ 0.0042 \\
3.007 & 0.9972 $\pm$ 0.0062 \\
3.698 & 0.9883 $\pm$ 0.0071 \\
\bottomrule\end{longtable}
In this emulator, v\_kin and eps\_kin therefore drive almost all the suppression up to k $\approx$ 3.7 h/Mpc, even with e\_w at its maximum.

\textbf{fiducial versus strongAGN.} The S(k) values are:

\begin{longtable}{@{}p{0.240\textwidth}p{0.240\textwidth}p{0.240\textwidth}p{0.240\textwidth}@{}}
\toprule
\textbf{Point} & \textbf{k = 1.01} & \textbf{k = 3.01} & \textbf{k = 3.70} \\ \midrule \endhead
fiducial & 0.9779 & 0.9197 & 0.8984 (std 0.004) \\
strongAGN & 0.9746 & 0.9121 & 0.8909 \\
\bottomrule\end{longtable}
Raising the kinetic parameters and e\_w at kappa\_w = 3 therefore adds only about 0.3--0.8 percentage points of suppression over this range. These values come from an earlier run on this thread and were not independently re-verified in the present run.

\textbf{strongAGN beyond k = 3.7 h/Mpc.} S(k) continues to fall:

\begin{longtable}{@{}p{0.480\textwidth}p{0.480\textwidth}@{}}
\toprule
\textbf{k (h/Mpc)} & \textbf{S(k)} \\ \midrule \endhead
5.25 & 0.851 \\
7.01 & 0.8205 \\
10.01 & 0.7946 (std 0.0098) \\
12.55 & 0.7893 \\
\bottomrule\end{longtable}
Between about 11.6 and 12.6 h/Mpc it flattens to a plateau of about 0.789--0.790, falling by less than 0.001 per 0.1 h/Mpc at the edge.

Two cautions apply to the high-k behaviour:

\begin{itemize}
  \item The ratio has no upturn within the grid. The minimum is at or beyond k\_max $\approx$ 12.6 h/Mpc, with a depth of about 21\% for strongAGN.
  \item maxsupp is already at 0.7356 by k = 7.84, deeper than the strongAGN plateau. Its minimum is unresolved.
\end{itemize}
The k\_max and plateau values come from byte offsets, and the final rows of the file were not quoted verbatim. They should be treated as approximate.

\paragraph{Redshift evolution, z = 0 to z = 1}
For maxsupp the suppression weakens strongly towards z = 1:

\begin{longtable}{@{}p{0.240\textwidth}p{0.240\textwidth}p{0.240\textwidth}p{0.240\textwidth}@{}}
\toprule
\textbf{k (h/Mpc)} & \textbf{S at z = 0} & \textbf{S at z = 1} & \textbf{Suppression, z = 0 to z = 1} \\ \midrule \endhead
1.013 & 0.9638 $\pm$ 0.0020 & 0.9781 $\pm$ 0.0018 & 3.6\% to 2.2\% \\
3.007 & 0.8698 $\pm$ 0.0065 & 0.9299 $\pm$ 0.0055 & 13.0\% to 7.0\% \\
3.70 & 0.8400 $\pm$ 0.0076 & 0.9214 $\pm$ 0.0064 & 16.0\% to 7.9\% \\
\bottomrule\end{longtable}
Qualitatively, this matches the FLAMINGO statement that the response grows as z falls (arXiv:2410.17109). That comparison has not been checked against the FLAMINGO text. The ratio is still falling at k = 3.7 at both redshifts, so any shift of the minimum with redshift could not be measured.

[[FIGURE 28]]

Figure 28 plots the emulator Pk-ratio at z = 0 and z = 1 for the fiducial and maxsupp points. Only the maxsupp values in the table above were read out numerically. The fiducial z = 1 prediction was written but not read, so no quantitative z-trend is claimed for the fiducial point.

A planned companion plot of the four z = 0 variants was overwritten by this one, because the two plotting calls shared an output path.

\paragraph{Gas fraction versus suppression}
The fGas emulator gives M\_gas/M500c inside R500c at z = 0 in three M500c bins: 5.01e13, 7.94e13 and 1.26e14 Msun/h.

\begin{longtable}{@{}p{0.320\textwidth}p{0.320\textwidth}p{0.320\textwidth}@{}}
\toprule
\textbf{Point} & \textbf{f\_gas at 1.26e14 Msun/h} & \textbf{f\_gas at 7.94e13 Msun/h} \\ \midrule \endhead
fiducial & 0.1214 $\pm$ 0.0010 & 0.1090 \\
strongAGN & 0.1221 $\pm$ 0.0019 & 0.1098 \\
weakAGN & 0.1414 $\pm$ 0.0023 & 0.1375 \\
\bottomrule\end{longtable}
Only fiducial and strongAGN have S(k) values to pair with these gas fractions. Their f\_gas values agree within errors, with strongAGN nominally higher, yet strongAGN is 0.3--0.8 points more suppressed. This pair therefore shows no resolvable van Daalen-type relation (arXiv:1906.00968), and the nominal trend runs the wrong way.

The comparison has two weaknesses:

\begin{itemize}
  \item It is confounded, because e\_w changes together with the kinetic parameters.
  \item It is limited, because the weakAGN S(k) was never read. weakAGN has a markedly higher gas fraction, so it is the point that would make a three-point test meaningful.
\end{itemize}
One physical reading of the pair is that the emulator's suppression is driven by gas removed beyond R500c, or by haloes of other masses, rather than by the R500c gas fraction. This is untested.

\paragraph{Comparison with observational suppression estimates}
Errors are combined in quadrature, taking the published error on the side facing the emulator value. For maxsupp:

\begin{longtable}{@{}p{0.480\textwidth}p{0.480\textwidth}@{}}
\toprule
\textbf{Constraint} & \textbf{Gap} \\ \midrule \endhead
Grandis et al., k = 1 & 0.4 sigma (4.2\% vs 3.6\%) \\
Grandis et al., k = 3 & 1.7 sigma (18.4\% vs 13.0\%) \\
eROSITA + kSZ, k = 1 & 3.2 sigma short (10\% vs 3.6\%) \\
HSC-XXL, k = 1 & 0.35 sigma \\
\bottomrule\end{longtable}
An earlier pass quoted about 3 and 8 sigma against Grandis et al. Those figures used emulator error alone and are superseded.

The strongest point, maxsupp\_lowM, is:

\begin{itemize}
  \item about 2.8 sigma short of eROSITA + kSZ;
  \item consistent with HSC-XXL and with Grandis et al. at k = 1;
  \item 3.6 points below Grandis et al. at k = 3.
\end{itemize}
All of these significances depend on literature values that have not been confirmed verbatim. They are also tested only at the kappa\_w = 2, e\_w = 1 corners with fixed cosmology.

\paragraph{Weak-lensing impact (literature-based)}
\textbf{Omega\_m bias from ignoring or mis-modelling baryons.} Maraio, Hall and Taylor (arXiv:2410.12500, Table 3, read verbatim) consider an idealised Euclid-like Stage-IV cosmic-shear survey with ell\_max = 5000:

\begin{longtable}{@{}p{0.240\textwidth}p{0.240\textwidth}p{0.240\textwidth}p{0.240\textwidth}@{}}
\toprule
\textbf{Ground truth} & \textbf{Model} & \textbf{100 $\Delta$$\Omega$\_m} & \textbf{Significance} \\ \midrule \endhead
Illustris & one-parameter HMCode, no theory error & -1.317 & -5.69 sigma \\
Illustris & same, with theoretical-error covariance & 0.128 & 0.249 sigma \\
Bahamas & one-parameter HMCode & 0.057 & 0.258 sigma \\
Eagle & six-parameter model, with theory error & 0.573 & 1.459 sigma \\
\bottomrule\end{longtable}
The text adds two statements:

\begin{itemize}
  \item The one-parameter model without theory error gives ''at least a 3 sigma offset in either Omega\_m or S8 for all hydro-sims excluding Bahamas''.
  \item With Simba as ground truth, even ell\_max = 1000 is biased by more than 2 sigma in all parameters.
\end{itemize}
The $\Delta$S8 values are in that paper's Table 4, which was not read. This bias is therefore an Omega\_m bias, not an S8 bias.

\textbf{Suppression needed to reconcile lensing with Planck.}

\begin{itemize}
  \item For DES Y3, A\_mod = 0.858 $\pm$ 0.052, with no scale cuts and a Planck prior (arXiv:2305.09827, Eq. 3). The suppression must extend to k \textasciitilde{} 0.2 h/Mpc.
  \item For KiDS-1000, A\_mod $\approx$ 0.69 (arXiv:2206.11794).
  \item Without A\_mod, DES Y3 S8 lies 0.9--1.8 sigma below Planck S8 = 0.828 $\pm$ 0.016, depending on analysis variant.
\end{itemize}
\textbf{Comparison with the emulator.} For maxsupp, S(k) is 0.9994 at k = 0.10 and 0.987 at k = 0.51 h/Mpc. The value at k $\approx$ 0.2 h/Mpc was not read but presumably lies below a 1\% effect. For maxsupp\_lowM the read value at k = 0.212 is 0.9967.

Emulator-level feedback therefore agrees with the Grandis et al. bound of less than 1\% for k \textless{} 0.37 h/Mpc. It appears far weaker than A\_mod = 0.86 or 0.69 would require at these scales. This is a qualitative comparison, because no A\_mod-equivalent number was computed for any emulator point.

\paragraph{Bibliographic and data products}
Validated citation entries for arXiv:2305.09827, arXiv:2410.12500, arXiv:2512.02954 and arXiv:2309.02920, each with a verbatim abstract quote, were written to \texttt{runs/\allowbreak{}261023\_\allowbreak{}143127/\allowbreak{}references.bib}.

Full texts of FLAMINGO (arXiv:2410.17109), Siegel et al. (arXiv:2512.02954) and Grandis et al. (arXiv:2309.02920) were downloaded to \texttt{runs/\allowbreak{}261023\_\allowbreak{}143127/\allowbreak{}} as \texttt{.txt} files. They have not yet been read.

\subsubsection*{Discussion}
\textbf{Feedback physics in the emulator.} Within this emulator, the kinetic AGN mode is the lever on S(k). Switching it off removes almost all suppression even at maximum e\_w. Lowering the seed mass, and plausibly lowering kappa\_w (not isolated, because M\_seed also changed), deepens it.

This is consistent with the general picture of AGN-driven gas expulsion from group-scale haloes. However, the emulator does not show the expected anti-correlation between the gas fraction inside R500c at about 1e14 Msun/h and the suppression. Possible reasons:

\begin{itemize}
  \item the relevant gas is pushed beyond R500c;
  \item other halo masses dominate;
  \item the two-point comparison is simply too confounded to show it.
\end{itemize}
The literature (arXiv:1906.00968, arXiv:2509.04552, arXiv:2309.06479) predicts a clear relation, so this is a real tension worth resolving with the unread weakAGN curve and single-parameter sweeps.

\textbf{Size of the suppression.} The emulator's largest sampled suppression at k = 1 h/Mpc, 4.4\%, agrees with HSC-XXL and Grandis et al. It falls short of the strong-feedback eROSITA + kSZ inference. At k = 3 h/Mpc it is below Grandis et al.

Three explanations are open:

\begin{itemize}
  \item the design range may not span the strongest feedback now favoured observationally;
  \item the cosmology axes, untouched here, may shift the response through the baryon fraction and halo abundance;
  \item the P\_grav reference may differ from P\_DMO.
\end{itemize}
\textbf{Weak lensing.} The literature

\end{document}
